% L002 student edition. Compile with xelatex twice (or latexmk -xelatex).
% All original reading passages, exercises, answer keys and bilingual supplements are included.
% This is a learner-facing document; audit logs and repository metadata are separate.
\documentclass[a4paper,11pt,openany]{report}
\usepackage{xeCJK}
\usepackage{fontspec}
\usepackage{geometry}
\geometry{a4paper,top=22mm,bottom=23mm,left=24.5mm,right=24.5mm,headheight=23pt,headsep=9mm,footskip=12mm}
\usepackage{microtype}
\usepackage{setspace}
\setstretch{1.20}
\usepackage{xcolor}
\definecolor{deepforest}{HTML}{20483F}
\definecolor{leaf}{HTML}{397665}
\definecolor{accentgold}{HTML}{B68A54}
\definecolor{softgreen}{HTML}{F4F8F5}
\definecolor{softsand}{HTML}{F9F7F1}
\definecolor{linegrey}{HTML}{DFE5DE}
\definecolor{bodygrey}{HTML}{4B5753}
\setmainfont{Latin Modern Roman}
\setsansfont{Latin Modern Sans}
\setmonofont{Latin Modern Mono}[Scale=0.90]
\setCJKmainfont[ItalicFont={Noto Serif CJK SC}]{Noto Serif CJK SC}
\setCJKsansfont{Noto Sans CJK SC}
\setCJKmonofont{Noto Sans CJK SC}
\newCJKfontfamily\CJKsf{Noto Sans CJK SC}
\usepackage[most]{tcolorbox}
\usepackage{enumitem}
\usepackage{needspace}
\usepackage{etoolbox}
\usepackage{ragged2e}
\usepackage{fancyhdr}
\usepackage{lastpage}
\usepackage{titlesec}
\usepackage{xurl}
\usepackage{longtable,booktabs,array}
\usepackage{hyperref}
\hypersetup{colorlinks=true,linkcolor=deepforest,urlcolor=leaf,bookmarksnumbered=false,pdfcreator={XeLaTeX},pdftitle={考研英语（一）词汇与阅读讲义：环境与生态},pdfsubject={L002 环境与生态：保护政策与野火管理}}
\urlstyle{same}
\sloppy
\emergencystretch=2em
\setlength{\parindent}{1.8em}
\setlength{\parskip}{0.14em}
\setcounter{secnumdepth}{-1}
\setcounter{tocdepth}{1}
\raggedbottom
\renewcommand{\contentsname}{学习目录}
\setlist[itemize]{leftmargin=2.0em,topsep=3pt,itemsep=3pt,parsep=0pt}
\setlist[enumerate]{leftmargin=2.0em,topsep=3pt,itemsep=3pt,parsep=0pt}
\setlist[description]{leftmargin=2.0em,topsep=3pt,itemsep=3pt}
% Titles: unit heading, subsection, and sentence drills all share a clear rhythm.
\titleformat{\chapter}[display]
 {\CJKsf\color{deepforest}\bfseries}
 {\small\color{accentgold} 考研英语（一）\, · \,第二册}
 {5pt}
 {\Huge}
 [\vspace{1.0ex}{\color{accentgold}\rule{\linewidth}{1.1pt}}]
\titlespacing*{\chapter}{0pt}{-2mm}{9mm}
\titleformat{\section}[hang]{\CJKsf\Large\bfseries\color{deepforest}}{}{0pt}{}
\titlespacing*{\section}{0pt}{5mm}{2mm}
\titleformat{\subsection}[hang]{\CJKsf\normalsize\bfseries\color{leaf}}{}{0pt}{}
\titlespacing*{\subsection}{0pt}{3mm}{0.8mm}
\titleformat{\subsubsection}[hang]{\CJKsf\normalsize\color{deepforest}}{}{0pt}{}
\titlespacing*{\subsubsection}{0pt}{3mm}{1mm}
\renewcommand{\chaptermark}[1]{\markboth{#1}{}}
\pagestyle{fancy}
\fancyhf{}
\fancyhead[L]{\CJKsf\footnotesize\color{leaf} 英语（一）\, · \,环境与生态}
\fancyhead[R]{\CJKsf\footnotesize\color{bodygrey}\nouppercase{\leftmark}}
\fancyfoot[C]{\color{bodygrey}\sffamily\small\thepage}
\renewcommand{\headrulewidth}{0.3pt}
\renewcommand{\headrule}{\hbox to\headwidth{\color{linegrey}\leaders\hrule height \headrulewidth\hfill}}
\fancypagestyle{plain}{\fancyhf{}\fancyfoot[C]{\sffamily\small\color{bodygrey}\thepage}\renewcommand{\headrulewidth}{0pt}}
% Breakable textbook word profile with compact labels, designed to sit across pages gracefully.
\newtcolorbox{vocabbox}[3]{enhanced,
 before skip=6pt,after skip=6pt,
 colback=softgreen,colframe=linegrey,boxrule=0.5pt,
 borderline west={2.3pt}{0pt}{leaf},arc=1.6mm,
 left=3mm,right=3mm,top=2.8mm,bottom=2.8mm,
 fonttitle=\CJKsf\bfseries\color{deepforest},colbacktitle=softgreen,
 attach boxed title to top left={xshift=3mm,yshift=-1mm},
 boxed title style={boxrule=0pt,frame hidden,colback=softgreen},
 title={\color{accentgold}\textsf{#1}\quad\large\textsf{#2}\quad\normalsize\color{bodygrey}\textsf{#3}},
 before upper={\parindent0pt\parskip=0.17em\relax},
 }
\newcommand{\emphasislabel}[1]{\textcolor{leaf}{\CJKsf\bfseries #1}}
\newcommand{\zhline}[1]{{\CJKsf\color{bodygrey}#1}}
% Block quotation styles: enough contrast, without repeatedly boxing the bilingual appendix.
\renewenvironment{quote}{%
 \par\begin{list}{}{\setlength{\leftmargin}{1.25em}\setlength{\rightmargin}{0em}}\item[]\color{bodygrey}\ignorespaces
}{\end{list}}
\providecommand{\tightlist}{\setlength{\itemsep}{1pt}\setlength{\parskip}{0pt}}
\providecommand{\pandocbounded}[1]{#1}
\providecommand{\hyperlink}[2]{\hyperref[#1]{#2}}
\providecommand{\hypertarget}[2]{\phantomsection\label{#1}#2}

\begin{document}
\begin{titlepage}
\thispagestyle{empty}
\sffamily
\vspace*{12mm}
{\CJKsf\color{leaf}\large 考研英语（一） \textbar{} 词汇与阅读讲义}\par
\vspace{40mm}
{\color{accentgold}\rule{22mm}{2pt}}\par
\vspace{10mm}
{\CJKsf\color{deepforest}\fontsize{34}{44}\selectfont\bfseries 环境与生态}\par
\vspace{4mm}
{\CJKsf\color{deepforest}\fontsize{20}{28}\selectfont\bfseries 保护政策与野火管理}\par
\vspace{13mm}
{\color{leaf}\fontsize{13}{18}\selectfont ENVIRONMENT \& ECOLOGY}\par
\vspace{3mm}
{\color{bodygrey}\fontsize{11}{17}\selectfont Wetlands \textperiodcentered{} Wildfire Management \textperiodcentered{} Critical Reading}\par
\vfill
\begin{tcolorbox}[colback=softsand,colframe=softsand,boxrule=0pt,arc=1mm,left=4mm,right=4mm,top=5mm,bottom=5mm]
{\CJKsf\color{deepforest}\bfseries 第二册 · 阅读与词汇训练}\par\vspace{1mm}
{\CJKsf\color{bodygrey}\normalsize 两篇真题伴读　｜　23 项重点词与词块　｜　三篇原创阅读}\par
{\CJKsf\color{bodygrey}\normalsize 句法和翻译练习　｜　答案详解　｜　英中双语对照}\par
\end{tcolorbox}
\vspace{12mm}
\end{titlepage}
\pagenumbering{roman}
\tableofcontents
\clearpage
\pagenumbering{arabic}
\chapter{导读｜开始学习之前}\label{ux5bfcux8bfbux5f00ux59cbux5b66ux4e60ux4e4bux524d}

这册有两个看似接近、实际上很不相同的故事。一篇谈美国切萨皮克湾湿地。河流与污染物可以跨越州界，但环境监管权力受到法律和行政边界的限制。另一篇谈森林野火。长期扑火的确保护过人和财产，但研究者发现，抑制所有小火也可能给未来留下更多可燃物。

本册借助环境议题训练英语阅读。学习时请关注：\textbf{一个动词究竟改变了什么？一个转折推翻的是事实还是先前的判断？作者又凭什么从事实走向结论？}

建议分四次学习：第一轮独立完成诊断及两篇真题（先打开试卷，再读下文导读）；第二轮精读前11项词汇和2024篇；第三轮精读后12项词汇和2026篇；第四轮完成原创阅读、专项练习、答案与双语核对。每次遇到不认识的词，先依据上下文尝试判断，再使用词汇部分；不要从第一页一直看到答案。

\textbf{真题阅读入口}（非官方试卷整理与第三方转录）：2024英语（一）Text 4：\href{https://ac-wiki.org/general-skills/study/assets/2024-201.pdf}{试卷 PDF}（PDF 第9---10页）。2026英语（一）Text 4：\href{https://english-exam.lazynote.cn/kaoyan/sections/2026-english-one/section2-part-a-4/}{在线题目与原文转录}（第三方转录）。这册提供真题导读与题目解析，不重印两篇全文；阅读前尽量先打开原卷独立作答。

\begin{quote}
阅读本册时，请先在纸上留下自己的答案：这有助于分辨``看懂了解析''和``自己真正会做''之间的差别。
\end{quote}

\clearpage
\section{开始前的4道小测}\label{ux5f00ux59cbux524dux76844ux9053ux5c0fux6d4b}

不查词典作答。这里只测试你对范围、冲突、程度与因果的敏感度；正确答案放在最后。

\Needspace{10\baselineskip}
\textbf{D01.} The ruling narrowed the scope of federal oversight.

\emph{What changed?}

A. The physical size of the wetlands.

B. The range of federal oversight.

C. The amount of rainfall.

D. The number of state governments.

\Needspace{10\baselineskip}
\textbf{D02.} The proposal was at odds with the findings of earlier studies.

\emph{The relationship between the proposal and the evidence is one of \_\_\_\_\_.}

A. agreement

B. uncertainty

C. conflict

D. repetition

\Needspace{10\baselineskip}
\textbf{D03.} Measures that curb pollution do not necessarily remove every source of it.

\emph{Which statement follows?}

A. The measures may reduce but not eliminate pollution.

B. The measures have no effect on pollution.

C. All pollution sources have disappeared.

D. The measures always increase pollution.

\Needspace{10\baselineskip}
\textbf{D04.} A policy meant to reduce short-term harm may create risks that emerge only later.

\emph{What does the sentence emphasize?}

A. The absence of policy benefits.

B. The certainty of policy failure.

C. The difficulty of passing new laws.

D. The possible tension between immediate and later effects.

\chapter{第一单元｜湿地：法律边界与水的边界}\label{ux7b2cux4e00ux5355ux5143ux6e7fux5730ux6cd5ux5f8bux8fb9ux754cux4e0eux6c34ux7684ux8fb9ux754c}

\section{01 为什么要读这篇文章}\label{ux4e3aux4ec0ux4e48ux8981ux8bfbux8fd9ux7bc7ux6587ux7ae0}

想象一条河流，流经几个行政区域。上游的农田施肥，下游的社区取水；一个地方批准开发，另一个地方可能承担径流与洪水的代价。环境问题并不总停留在作出决定的地方。这是理解2024年英语（一）Text 4的第一把钥匙。

切萨皮克湾（Chesapeake Bay）位于美国东部。文章开头介绍湿地的多重功能：提供生境、蓄纳洪水、过滤污染物，并降低风暴潮带来的风险。随后引入一项法院裁决，说明联邦环境保护局（EPA）能够监管的湿地范围受到收缩。作者关心的并非抽象的``政府是否应该干预''，而是\textbf{水体的生态联系如何面对法律管辖边界}。

第一次读时，先圈出这些角色：\textbf{湿地---最高法院---EPA---各州---上游农民---下游社区---州立法者}。再用箭头画出它们之间的关系。

\section{02 先读真题：只回答这五个问题}\label{ux5148ux8bfbux771fux9898ux53eaux56deux7b54ux8fd9ux4e94ux4e2aux95eeux9898}

打开2024真题 Text 4，独立完成第36---40题，再回到这里写简短答案。本册以五个论证环节定位；部分网页把同一环节拆成数段，须同时核对内容关键词。

\begin{itemize}
\tightlist
\item
  \textbf{第36题}：首段提供了哪些具体事实，使``这一自然环境有价值''的评价成立？至少找出两项。
\item
  \textbf{第37题}：裁决改变的是湿地面积、污染量，还是EPA行使权力的范围？
\item
  \textbf{第38题}：作者担忧的是``已经发生的全部后果''，还是``可能出现的未来问题''？用原文评价语判断。
\item
  \textbf{第39题}：为什么一个跨州流域需要协调保护？哪些因果环节连接不同地区？
\item
  \textbf{第40题}：结尾的行动号召面向什么人？作者要求他们重新审视什么？
\end{itemize}

别急着背答案。阅读训练真正有效的时刻，是你必须向自己解释``为什么文本支持这一项，而不支持另一项''。

\section{03 按五个论证环节精读}\label{ux6309ux4e94ux4e2aux8bbaux8bc1ux73afux8282ux7cbeux8bfb}

以下P1---P5是本册伴读定位；网页分段可能不同，以试卷和内容锚点共同核对。

\begin{itemize}
\tightlist
\item
  \textbf{第1段｜建立价值}　湿地通过生态服务说明其价值，而非单靠商业收益。
\item
  \textbf{第2段｜政策变化}　法院裁决缩小联邦监管范围，先弄清改变了什么。
\item
  \textbf{第3段｜质疑局部视角}　州内存在保护法，并不意味着跨州问题已经解决。
\item
  \textbf{第4段｜解释协调必要性}　上游决策可以影响下游，联邦参与有助于跨地区协调。
\item
  \textbf{第5段｜提出行动请求}　作者呼吁州立法者重视保护；呼吁不等于行动已经发生。
\end{itemize}

\emphasislabel{最容易读错的地方}：第2段不是说``湿地已经完全失去保护''，第3段也不是说``所有州都会采取同样的保护''；作者讨论的是监管权限收缩带来的风险和协调问题。

\subsection{精读示范：如何从``本州有法律''走到``仍不能放心''}\label{ux7cbeux8bfbux793aux8303ux5982ux4f55ux4eceux672cux5ddeux6709ux6cd5ux5f8bux8d70ux5230ux4ecdux4e0dux80fdux653eux5fc3}

先写下看似令人安心的前提：本州已有湿地保护法。这个前提本身没有被作者否定；作者否定的是从此前提直接推出``本州水域已足够安全''。缺失的中间条件是：影响本州水质的行为是否全部发生在本州？流域跨界的事实表明并非如此。

因此，读到 shortsighted，不能只把它翻译成``短视''。应找出它评价的对象------仅凭本州法律就感到安心的看法；再找理由------外部辖区的决定会沿着水流影响本州。随后出现上游施肥的例子，是把抽象的协调需求具体化，并不是换了一个农业主题。

\textbf{停下来作答（1）}　用``前提---遗漏条件---结论''写三行笔记。把州法律存在和跨州保护充分分别放入哪一行？

\textbf{停下来作答（2）}　如果上游利益与下游利益并不一致，为什么``各方只管自己的区域''不足以回应作者的担忧？答案应写出传播路径，不只写``需要合作''。

\textbf{如何识别选项改写}　先把证据还原成关系：跨界影响使协调不可少。正确概括可以不重复农田、蓝蟹等细节，但必须保留关系；错误概括则可能把协调需要写成已经解决所有利益冲突。概括允许删去例子，不允许添加新成效。

\section{04 词汇精读｜湿地与监管}\label{ux8bcdux6c47ux7cbeux8bfbux6e7fux5730ux4e0eux76d1ux7ba1}

新词不按字母排列，而按文章中实际出现的功能顺序学习：\textbf{环境作用 \(\rightarrow\) 监管边界 \(\rightarrow\) 后果评价 \(\rightarrow\) 行动号召}。每个词有两句原创英文材料。先读英文并尝试口头翻译，完成本册之后再到双语附录核对。

\begin{vocabbox}{01}{filter}{v.}

\textbf{这篇文章中的意思}　过滤；滤除

\textbf{在真题里}　2024·伴读环节P1｜湿地净化水体

\textbf{常用搭配}　filter pollutants · filter out impurities

\textbf{义项辨析}　filter out 强调将杂质滤除；filter information 也可用于信息筛选

\textbf{例句 01-1}　Wetland vegetation can filter some pollutants before they reach open water.

\textbf{例句 01-2}　The committee filtered out unreliable reports before reviewing the evidence.

\emphasislabel{读题提醒}　先确认被过滤的对象。文章说湿地过滤污染物，并不是让污染物``消失于所有地方''。

\textbf{用法再进一层}　filter 是及物动词，直接接被处理的物质；filter out 的宾语是被剔除的部分。看到 out，要判断谁被留下、谁被去掉。

\end{vocabbox}

\begin{vocabbox}{02}{authority}{n.}

\textbf{这篇文章中的意思}　法定监管权限（不可数）；其他语境中也可指当局或权威人士

\textbf{在真题里}　2024·伴读环节P2｜环保部门可监管的范围

\textbf{常用搭配}　legal authority · regulatory authority · local authorities

\textbf{义项辨析}　authority 强调有没有权力去做；jurisdiction 更常指权力所适用的地域或事务范围

\textbf{例句 02-1}　The agency has the authority to inspect factories that discharge waste into the river.

\textbf{例句 02-2}　Several local authorities expressed concern about the proposed changes.

\emphasislabel{读题提醒}　本篇的 authority 是监管权限，不是``某位科学权威''。

\textbf{用法再进一层}　表示权限时通常不可数：have authority to do；表示某一领域的权威人士时可数：an authority on a subject。复数 authorities 常指主管部门。

\end{vocabbox}

\begin{vocabbox}{03}{scope}{n.}

\textbf{这篇文章中的意思}　监管的适用范围；``余地''属于其他语境中的义项

\textbf{在真题里}　2024·伴读环节P2｜环保监管的适用边界

\textbf{常用搭配}　the scope of a law · beyond the scope of · narrow/broaden the scope

\textbf{义项辨析}　scope 与 scale 不同，前者重在涵盖内容，后者重在规模或大小

\textbf{例句 03-1}　The decision narrowed the scope of the regulation without abolishing it.

\textbf{例句 03-2}　The report is valuable, although several issues fall beyond its scope.

\emphasislabel{读题提醒}　注意``缩小适用范围''并不等于``完全取消监管''。

\textbf{用法再进一层}　within/beyond the scope of 后接某项研究、法规或讨论。scope for improvement 中 scope 指改进余地；不要套入本篇的监管范围义。

\end{vocabbox}

\begin{vocabbox}{04}{at odds with}{phr.}

\textbf{这篇文章中的意思}　与\ldots\ldots 不一致；相冲突

\textbf{在真题里}　2024·伴读环节P2｜商业利益与环保规则不一致

\textbf{常用搭配}　be at odds with evidence · at odds with public expectations

\textbf{义项辨析}　不等于 against the law（违法），它首先表示不协调或不一致

\textbf{例句 04-1}　The policy is at odds with the evidence gathered by independent researchers.

\textbf{例句 04-2}　Her public statement seemed at odds with the company's earlier promises.

\emphasislabel{读题提醒}　选择题往往把``不一致''扩大为``违法''或``完全敌对''，要保留原文的强度。

\textbf{用法再进一层}　完整结构是 A be at odds with B。with 后的对象可以是人、政策或证据；介词结构本身不判断谁正确。

\end{vocabbox}

\begin{vocabbox}{05}{repercussion}{n.}

\textbf{这篇文章中的意思}　连带后果；后续影响（常为负面）

\textbf{在真题里}　2024·伴读环节P2｜裁决可能造成的环境后果

\textbf{常用搭配}　have repercussions for · serious repercussions · unintended repercussions

\textbf{义项辨析}　consequence 可中性，repercussion 往往侧重不利的间接后果

\textbf{例句 05-1}　A change in upstream farming may have repercussions for communities downstream.

\textbf{例句 05-2}　The decision had unintended repercussions for small businesses.

\emphasislabel{读题提醒}　辨别``可能有重大后果''和``后果已经全部发生''。

\textbf{用法再进一层}　常以复数出现，have repercussions for 指对谁造成后续影响。原文情态词和时态决定后果是可能、已经发生还是尚未落实。

\end{vocabbox}

\begin{vocabbox}{06}{jurisdiction}{n.}

\textbf{这篇文章中的意思}　管辖权；管辖区域

\textbf{在真题里}　2024·伴读环节P3｜流域跨越多个州

\textbf{常用搭配}　within a jurisdiction · beyond federal jurisdiction · neighboring jurisdictions

\textbf{义项辨析}　在本篇中可用``权限''与``辖区''帮助区分；jurisdiction 本身也能表示管辖权

\textbf{例句 06-1}　The river crosses several jurisdictions, each with different environmental rules.

\textbf{例句 06-2}　The dispute concerns whether the case falls within federal jurisdiction.

\emphasislabel{读题提醒}　污染物越过边界，不代表法律管辖权会自动延伸。

\textbf{用法再进一层}　管辖权义常不可数，如 under federal jurisdiction；行政区域义可数，如 several jurisdictions。authority 也可以表示法定权限，两词有重叠，区别不是绝对互斥。

\end{vocabbox}

\begin{vocabbox}{07}{shortsighted}{adj.}

\textbf{这篇文章中的意思}　目光短浅的；缺乏长远考虑的

\textbf{在真题里}　2024·伴读环节P3｜作者对``只看本州保护''的评价

\textbf{常用搭配}　a shortsighted view · a shortsighted policy

\textbf{义项辨析}　常是评价性形容词；须确定是作者判断还是引述他人判断

\textbf{例句 07-1}　It would be shortsighted to ignore maintenance costs when choosing a cheaper system.

\textbf{例句 07-2}　A shortsighted decision may bring savings now but greater costs later.

\emphasislabel{读题提醒}　读到评价词时，应追问``作者在批评哪个看法''。

\textbf{用法再进一层}　shortsighted 批评的是判断未充分考虑更远的后果。原文也可能强调超出本地区的空间影响，不应只理解为时间上的短期。

\end{vocabbox}

\begin{vocabbox}{08}{transcend}{v.}

\textbf{这篇文章中的意思}　超越；跨越（限制、边界或分歧）

\textbf{在真题里}　2024·伴读环节P4｜跨越州界及特殊利益的影响

\textbf{常用搭配}　transcend borders · transcend political divisions

\textbf{义项辨析}　transcend 在议论文中常接抽象边界、局限、差异

\textbf{例句 08-1}　Effective river management must transcend administrative boundaries.

\textbf{例句 08-2}　The exhibition aims to transcend differences between generations.

\emphasislabel{读题提醒}　不要把它仅译成空间意义的``穿过''。

\textbf{用法再进一层}　transcend 直接接宾语，不加 across：transcend boundaries。它表示超出局限，不等于已经解决边界两边的全部冲突。

\end{vocabbox}

\begin{vocabbox}{09}{runoff}{n.}

\textbf{这篇文章中的意思}　地表径流，即在地表流动、汇入河湖的水

\textbf{在真题里}　2024·伴读环节P4｜污染物由上游传向下游

\textbf{常用搭配}　agricultural runoff · stormwater runoff · runoff into streams

\textbf{义项辨析}　runoff 是运动介质，不等于全部污染物本身

\textbf{例句 09-1}　Heavy rain generated runoff that carried soil from fields into nearby streams.

\textbf{例句 09-2}　Vegetated buffer strips can filter pollutants from runoff before it reaches the lake.

\emphasislabel{读题提醒}　阅读时留意``上游决定---径流---下游影响''的因果链。

\textbf{用法再进一层}　环境语境中 runoff 通常不可数。run off 可作两个词的动词短语；阅读时要区分名词与动词短语。

\end{vocabbox}

\begin{vocabbox}{10}{call on}{phr.}

\textbf{这篇文章中的意思}　呼吁；敦促；（另义）拜访

\textbf{在真题里}　2024·伴读环节P5｜作者向州立法者提出请求

\textbf{常用搭配}　call on lawmakers to act · call on the public to support

\textbf{义项辨析}　call on sb. to do 是呼吁某人做事；call on sb. 无后接不定式时也可意为拜访

\textbf{例句 10-1}　The editorial called on lawmakers to review existing wetland protections.

\textbf{例句 10-2}　The researchers called on schools to reconsider how success is measured.

\emphasislabel{读题提醒}　特别关注 sb. to do 的结构，找清行动主体与行动内容。

\textbf{用法再进一层}　call on sb. to do 中 sb. 是被呼吁采取行动的人。called on officials to act 中真正要行动的是 officials，不是 editorial。

\end{vocabbox}

\begin{vocabbox}{11}{stakes}{n. pl.}

\textbf{这篇文章中的意思}　利害关系；相关的重要利益

\textbf{在真题里}　2024·伴读环节P5｜保护湿地所牵涉的重大利害

\textbf{常用搭配}　the stakes are high · what is at stake · raise the stakes

\textbf{义项辨析}　不要把 stakes 机械译作博彩``赌注''；at stake 是``处于风险之中''

\textbf{例句 11-1}　The stakes are high for villages that depend on the river for drinking water.

\textbf{例句 11-2}　What is at stake is not just a single species but the stability of the wider ecosystem.

\emphasislabel{读题提醒}　可借助上下文理解``谁可能受益或受损''。

\textbf{用法再进一层}　the stakes are high 用复数 stakes；at stake 是固定形式，不写 at stakes。两者均要追问牵涉谁的利益。

\end{vocabbox}

\section{05 把词重新放回文章}\label{ux628aux8bcdux91cdux65b0ux653eux56deux6587ux7ae0}

重点不只是说出 \texttt{authority} 和 \texttt{jurisdiction} 的中文。请用它们回答以下问题。

\textbf{问题一}　为什么``EPA的 authority 受到限制''和``水中污染减少''不是同一件事？前者是什么类型的变化，后者又需要什么证据？

\textbf{问题二}　\texttt{at\ odds\ with} 表示不一致。假如选项说``开发企业已经违法''，它比原文多出了什么判断？

\textbf{问题三}　\texttt{repercussions} 表示连带后果。为什么作者需要从法院裁决谈到跨州径流，才能说明这些后果？

\textbf{问题四}　作者用 \texttt{shortsighted} 批评的是哪一种想法？这只是情绪表达，还是能够从流域事实中获得支持？

\begin{quote}
做完后用一句话概括：\textbf{湿地价值（事实）\(\rightarrow\) 监管变化（事件）\(\rightarrow\) 跨州后果（推理）\(\rightarrow\) 审查保护（主张）}。
\end{quote}

\section{06 词组与近义词的真正区别}\label{ux8bcdux7ec4ux4e0eux8fd1ux4e49ux8bcdux7684ux771fux6b63ux533aux522b}

\begin{itemize}
\tightlist
\item
  \textbf{authority / jurisdiction}　两词都可以指法定权限；authority 在本篇侧重可行使的权力，jurisdictions 则指不同辖区。应由搭配和单复数判断，不能硬分成互不重叠的词义。
\item
  \textbf{scope}　规定涵盖的对象和情形，不是河流的物理面积。
\item
  \textbf{repercussion}　决策引起的后续连带影响；不保证全部后果已经发生。
\item
  \textbf{at odds with}　两种要求或立场冲突，但不等于行为必然违法。
\item
  \textbf{call on sb. to do}　明确``谁呼吁---谁行动---采取何种行动''；作者提出呼吁不等于行动已经完成。
\end{itemize}

回到原文把这几项所在的完整句子重新读一遍。若翻译仍不顺，首先检查介词、定语和指代，而不只是重复背词义。

\chapter{第二单元｜野火：为什么``火更少''不一定``更安全''}\label{ux7b2cux4e8cux5355ux5143ux91ceux706bux4e3aux4ec0ux4e48ux706bux66f4ux5c11ux4e0dux4e00ux5b9aux66f4ux5b89ux5168}

\section{01 从一个反直觉发现进入文章}\label{ux4eceux4e00ux4e2aux53cdux76f4ux89c9ux53d1ux73b0ux8fdbux5165ux6587ux7ae0}

如果一种政策让火灾发生次数明显减少，你会认为它一定成功吗？2026年英语（一）Text 4的前半篇恰好展示了这个判断为什么值得重新检查。

一位曾在美国林务局工作的研究者，在学习树轮研究后发现：一些古树记录了欧洲殖民者到来之前反复发生的火灾。进一步的研究显示，某些生态系统中的低烈度火灾可以清除地表枯枝、缓解树木竞争、抑制病害，并减少未来火灾可能使用的燃料。\textbf{``火灾次数''与``严重火灾的风险''不能被简单等同。}

文章最后并不是鼓励任由所有火燃烧。它介绍一种有计划、依赖天气等条件实施的管理方式，同时承认气候变化使这一工作更必要，也更困难。

\section{02 独立完成真题第36---40题}\label{ux72ecux7acbux5b8cux6210ux771fux9898ux7b2c3640ux9898}

打开2026年英语（一）Text 4。第一次阅读请做出以下标记：

\begin{itemize}
\tightlist
\item
  在第1段圈出研究者\textbf{感到意外的对象}，不要只记他最初承担的职责。
\item
  在第2段标出低烈度火灾的\textbf{至少三种作用}，试着把它们概括成一句话。
\item
  在第3段标出比较的\textbf{两组时间资料}，并写下``缺口''指什么。特别留意比较所适用的研究地点。
\item
  在第4段把``现代灭火的成效''和``现代灭火的副作用''用两种符号标记。二者可以同时成立。
\item
  在第5---6段分别写出Martin与Prichard的立场，再判断共同点。
\end{itemize}

最后选择36---40题选项。重点记录第38题中的\textbf{frequency}、第39题中 \textbf{may} 的作用，别因为一个词而扩大整段结论。

\section{03 六段论证结构}\label{ux516dux6bb5ux8bbaux8bc1ux7ed3ux6784}

\begin{itemize}
\tightlist
\item
  \textbf{第1段｜发现}　研究者从树轮火痕读出殖民前的火灾历史；要找准``惊讶的对象''。
\item
  \textbf{第2段｜解释}　展示低烈度火给森林带来的好处；不要把其中一种好处当段落主旨。
\item
  \textbf{第3段｜比较}　用历史记录与近代记录比较，发现研究点位的火灾次数缺口。
\item
  \textbf{第4段｜权衡}　现代灭火确实有效，长期抑制也可能增加未来灾害风险。
\item
  \textbf{第5段｜实践}　有计划的烧除在适宜条件下可以成为管理手段。
\item
  \textbf{第6段｜进一步限定}　气候变化使这种手段更有必要，也使实施更艰难。
\end{itemize}

需要时可把第4段画成两条线：一条是现代消防减少部分火灾，一条是长期积累可燃物使某些严重火灾更难控制。\textbf{两条线并不必然互相否定。}

\subsection{精读示范：三个比较对象不能混在一起}\label{ux7cbeux8bfbux793aux8303ux4e09ux4e2aux6bd4ux8f83ux5bf9ux8c61ux4e0dux80fdux6df7ux5728ux4e00ux8d77}

第3段先比较历史节律与近几十年的记录，再用``没有长期扑火时原本可能发生多少''的估计作为参照。实际记录与假设情形不是两次都已观察到的事实。读 shortfall 时，要同时标出研究点位、比较时段及参照值；不能只记``少了''，就把它扩展成世界各地野火都在减少。

第4段的转折也不是撤回前面的灭火成效。一个措施可以有效降低发生次数，又在更长时间内改变森林燃料状况。两种判断使用的评价对象不同。作者把气候变化一并纳入，因此不能把所有严重野火都归因于扑火，也不能说扑火毫无价值。

最后两段回答``接下来怎么办''。先读所提出的干预，再读实施条件。计划烧除具有潜在用途，但天气、人员和风险限制属于措施的一部分。脱掉条件后再复述结论，就会把有条件支持误读成无条件支持。

\textbf{停下来作答（3）}　假设某研究点位在一个相同时段记录20次火灾，而无扑火情形的估计值为100次。实际值是参照值的百分之几？差额是多少？这些数字能否告诉你损失减少了多少？本题数字为自编练习，不是真题数据。

\textbf{停下来作答（4）}　请分别写一句话说明现代灭火的已观察成效、可能的长期代价，以及计划烧除的条件。三句话不要互相替代。

\section{04 词汇精读｜发现、风险与管理}\label{ux8bcdux6c47ux7cbeux8bfbux53d1ux73b0ux98ceux9669ux4e0eux7ba1ux7406}

\begin{vocabbox}{12}{mandate}{n./v.}

\textbf{这篇文章中的意思}　被赋予的职责（名词）；动词 mandate 可表示强制规定

\textbf{在真题里}　2026·P1｜机构最初承担的灭火职责

\textbf{常用搭配}　a mandate to do · under its mandate · mandate changes

\textbf{义项辨析}　此处指赋予机构的职责，而不是所有语境都译作``强制命令''

\textbf{例句 12-1}　The forest service received a mandate to protect nearby communities from fire.

\textbf{例句 12-2}　A new law may mandate regular safety inspections for public buildings.

\emphasislabel{读题提醒}　机构接受的任务，不应误解成科学结论。

\textbf{用法再进一层}　a mandate to do 中 to do 说明职责内容；mandate inspections 中 mandate 是及物动词。名词与动词结构不可混用。

\end{vocabbox}

\begin{vocabbox}{13}{curb}{v./n.}

\textbf{这篇文章中的意思}　抑制；控制；遏制

\textbf{在真题里}　2026·P2｜减少病害传播

\textbf{常用搭配}　curb the spread of disease · curb emissions · curb spending

\textbf{义项辨析}　curb \(\neq\) eliminate：抑制不保证完全消除

\textbf{例句 13-1}　Small fires can help curb the spread of some forest diseases.

\textbf{例句 13-2}　The city introduced measures to curb excessive water use during the drought.

\emphasislabel{读题提醒}　题目容易把``减少风险''升级为``彻底杜绝''。

\textbf{用法再进一层}　curb 作动词直接接风险、排放或传播；a curb on spending 中 curb 是可数名词，指限制措施。降低程度不自动等于归零。

\end{vocabbox}

\begin{vocabbox}{14}{peter out}{phr.}

\textbf{这篇文章中的意思}　逐渐减弱并停止

\textbf{在真题里}　2026·P2｜低强度火自行减弱

\textbf{常用搭配}　the fire petered out · interest petered out

\textbf{义项辨析}　强调一个过程：由强到弱，最后终止

\textbf{例句 14-1}　The fire petered out when it reached an area with little dry vegetation.

\textbf{例句 14-2}　Public interest in the campaign petered out after several months.

\emphasislabel{读题提醒}　留意这是不及物短语，不说 \emph{peter out the fire}。

\textbf{用法再进一层}　peter out 不带宾语；过去式 petered out。不要把``慢慢结束''的过程义改写成 sudden disappearance。

\end{vocabbox}

\begin{vocabbox}{15}{compile}{v.}

\textbf{这篇文章中的意思}　收集并整理；汇编

\textbf{在真题里}　2026·P3｜汇集多年火灾记录

\textbf{常用搭配}　compile data · compile records · compile a report

\textbf{义项辨析}　compile 是将资料整理成系统的整体，不表示凭空制造证据

\textbf{例句 15-1}　The scientists compiled fire records from several research sites.

\textbf{例句 15-2}　The museum compiled visitors' comments into an annual report.

\emphasislabel{读题提醒}　既要看研究对象，也要看形成什么记录。

\textbf{用法再进一层}　compile A from B 表示从B汇集成A；compile A into B 表示把A整理成B。from 与 into 分别提示资料来源与成品。

\end{vocabbox}

\begin{vocabbox}{16}{shortfall}{n.}

\textbf{这篇文章中的意思}　短缺；不足；缺口

\textbf{在真题里}　2026·P3｜观测火灾次数少于历史预期

\textbf{常用搭配}　a shortfall in funding · a shortfall of supplies

\textbf{义项辨析}　本篇是数量相对于比较基准的不足，而非``灾害一定较少''

\textbf{例句 16-1}　The records revealed a shortfall in small fires compared with the number expected at the site.

\textbf{例句 16-2}　A funding shortfall forced the town to delay the restoration project.

\emphasislabel{读题提醒}　读到 shortfall 时，问清``实际值比哪个预期少''。

\textbf{用法再进一层}　a shortfall in funding 指哪方面不足；a shortfall of ten units 指具体缺多少。解释任何缺口，都应先说比较基准。

\end{vocabbox}

\begin{vocabbox}{17}{suppression}{n.}

\textbf{这篇文章中的意思}　压制；抑制；（火灾）扑灭

\textbf{在真题里}　2026·P3--P4｜长期灭火造成可燃物累积

\textbf{常用搭配}　fire suppression · suppress fires · suppression efforts

\textbf{义项辨析}　suppression 代表人为积极扑火，不代表所有环境管理措施

\textbf{例句 17-1}　Decades of fire suppression allowed dead wood to accumulate.

\textbf{例句 17-2}　The report examines whether strict suppression remains the best policy.

\emphasislabel{读题提醒}　与 prescribed burns（计划烧除）是不同治理手段。

\textbf{用法再进一层}　suppression 是 suppress 的名词形式。fire suppression 指扑火活动或政策，不能由这个名词单独推断所有火都成功扑灭。

\end{vocabbox}

\begin{vocabbox}{18}{testament}{n.}

\textbf{这篇文章中的意思}　体现、明证；``遗嘱''是其他语境中的义项

\textbf{在真题里}　2026·P4｜历史记录体现现代灭火的有效性

\textbf{常用搭配}　a testament to resilience · a testament to effort

\textbf{义项辨析}　a testament to 的 to 是介词，后接名词或动名词

\textbf{例句 18-1}　The survival of ancient trees is a testament to the forest's resilience.

\textbf{例句 18-2}　The new library is a testament to years of community cooperation.

\emphasislabel{读题提醒}　要先确认作者``以什么事实证明什么''，不要在段内割裂语义。

\textbf{用法再进一层}　a testament to 后面的 to 是介词。先找作为证据的事实，再找它所支持的评价；``体现有效性''不等于证明没有副作用。

\end{vocabbox}

\begin{vocabbox}{19}{erode}{v.}

\textbf{这篇文章中的意思}　侵蚀；逐渐削弱

\textbf{在真题里}　2026·P4｜极端天气下扑火能力受到削弱

\textbf{常用搭配}　erode confidence · erode trust · erode the ability to

\textbf{义项辨析}　erode \(\neq\) destroy：程度上可能未完全消失

\textbf{例句 19-1}　A growing fuel load can erode firefighters' ability to control severe fires.

\textbf{例句 19-2}　Repeated delays eroded public confidence in the transport plan.

\emphasislabel{读题提醒}　常见宾语是抽象事物，作者可能用比喻表达慢性恶化。

\textbf{用法再进一层}　erode 可及物也可不及物；本册关注 erode confidence/ability。逐渐削弱说明变化方向，实际损害程度仍须读上下文。

\end{vocabbox}

\begin{vocabbox}{20}{ignite}{v.}

\textbf{这篇文章中的意思}　起火（本句不及物）；也可及物，表示点燃或引发

\textbf{在真题里}　2026·P4｜恶劣天气条件下起火

\textbf{常用搭配}　ignite dry vegetation · ignite debate · ignite interest

\textbf{义项辨析}　同一动词搭配不同宾语，物理义与抽象义切换

\textbf{例句 20-1}　Strong winds can help a small spark ignite dry grass.

\textbf{例句 20-2}　The proposal ignited a debate over who should pay for flood protection.

\emphasislabel{读题提醒}　不要把 ignite debate 译成``点燃一场真实的火''。

\textbf{用法再进一层}　ignite 可不及物表示起火，也可及物表示点燃。天气可助长燃烧，但不要把干燥条件当成每场火的点火源。

\end{vocabbox}

\begin{vocabbox}{21}{prescribed}{adj.}

\textbf{这篇文章中的意思}　规定的；按计划进行的

\textbf{在真题里}　2026·P5｜计划烧除

\textbf{常用搭配}　prescribed burns · prescribed treatment · a prescribed method

\textbf{义项辨析}　prescribed burn 是专业固定搭配，指经过安排的、有控制条件的烧除

\textbf{例句 21-1}　Prescribed burns are conducted under carefully selected conditions.

\textbf{例句 21-2}　The prescribed procedure must be followed when dangerous materials are handled.

\emphasislabel{读题提醒}　它并不是``让大火自行燃烧''，更不是``开药方''。

\textbf{用法再进一层}　prescribed 是 prescribe 的过去分词，可作形容词。prescribed burns 是计划烧除，不保证实施中绝无风险。

\end{vocabbox}

\begin{vocabbox}{22}{tactic}{n.}

\textbf{这篇文章中的意思}　具体策略；手段

\textbf{在真题里}　2026·P6｜计划烧除作为治理手段

\textbf{常用搭配}　a practical tactic · a tactic for reducing risk

\textbf{义项辨析}　tactic 是具体做法；strategy 通常表示统筹性的长远战略

\textbf{例句 22-1}　Careful thinning is one tactic for reducing the amount of fuel near homes.

\textbf{例句 22-2}　The campaign adopted a different tactic after its early efforts failed.

\emphasislabel{读题提醒}　题目不应从``一种有效手段''推出``唯一政策''。

\textbf{用法再进一层}　one tactic for doing 表示实现目标的一种做法。one 保留其他可能手段，不能转述成 the only effective policy。

\end{vocabbox}

\begin{vocabbox}{23}{double-edged sword}{phr.}

\textbf{这篇文章中的意思}　双刃剑；本段借此比喻``更需要开展，却更难开展''的两难

\textbf{在真题里}　2026·P6｜气候变化使治理更必要但更困难

\textbf{常用搭配}　a double-edged sword · prove to be a double-edged sword

\textbf{义项辨析}　常见义是兼有利弊；此段不是说气候变化有益，而是说它同时增加治理需要与实施难度

\textbf{例句 23-1}　Greater public access can be a double-edged sword for nature reserves.

\textbf{例句 23-2}　Digital monitoring is a double-edged sword when privacy safeguards are weak.

\emphasislabel{读题提醒}　查明``什么方面需要更多''与``什么方面变得更难''。

\textbf{用法再进一层}　双刃剑的通常含义不能替代当前上下文。此篇的两面是需求增加与操作困难，不能编出``气候变化改善生态''的好处。

\end{vocabbox}

\section{05 跨册复习：三个旧词}\label{ux8de8ux518cux590dux4e60ux4e09ux4e2aux65e7ux8bcd}

这些词在前面的学习材料中已经出现，因此本册只做情境复现，不把它们当成全新重点词。

\begin{itemize}
\tightlist
\item
  \textbf{regulate}　监管、规范。原来在科技领域学习；现在用于辨别法律监管范围与环境实际变化。
\item
  \textbf{resilient}　有恢复力的、能够承受冲击并恢复。现在接入森林生态恢复。
\item
  \textbf{subsequent}　后续的、随后的。它提示结果发生的时间顺序，但不自动构成因果关系。
\end{itemize}

翻译下面三句复习材料，答案在最后。

\textbf{R01}　A resilient forest can recover from some disturbances, but its recovery depends on local conditions.

\textbf{R02}　Subsequent reports questioned whether the initial improvement would last.

\textbf{R03}　A government may regulate construction without having the power to prohibit every project.

\section{06 四组最容易混淆的表达}\label{ux56dbux7ec4ux6700ux5bb9ux6613ux6df7ux6dc6ux7684ux8868ux8fbe}

\textbf{curb / suppress / eliminate}　\texttt{curb} 强调遏制发展或传播；\texttt{suppress} 在本文指压制、扑灭；\texttt{eliminate} 才是消除。读到``长期 suppression''，要问压制了何种火，也要寻找可能产生的代价。

\textbf{frequency / severity / damage}　frequency 是发生频次，severity 是严重程度，damage 是造成的损失。三者有关联，但任何一个指标都不能替代全部指标。

\textbf{mandate / tactic}　mandate 是机构要履行的职责；tactic 是用来完成目标的一种具体方法。职责可以不变，手段却可以随着证据发生改变。

\textbf{prescribed burn / uncontrolled wildfire}　前者按条件和方案实施；后者并不处于相同控制之下。作者支持前者，不能据此推出``支持放任所有野火''。

\textbf{一句话收束}　真正需要读懂的是：一项措施为什么在某种意义上有效，又为什么在另一时间尺度或情景下出现局限。

\chapter{第三单元｜读懂句子里的转折、限定与证据}\label{ux7b2cux4e09ux5355ux5143ux8bfbux61c2ux53e5ux5b50ux91ccux7684ux8f6cux6298ux9650ux5b9aux4e0eux8bc1ux636e}

考研英语（一）的长句不一定难在语法名称，而常常难在\textbf{逻辑关系被套进多层修饰中}。下面8句全部原创，服务于这两篇文章的阅读问题。先划出主干，再用下方解析检查结构；译文留到末尾附录。

方法：先找主语和谓语，再标出 \texttt{although\ /\ but\ /\ that\ /\ whether\ /\ rather\ than\ /\ only\ after} 等连接词，最后判断句子的结论是否被 \texttt{may\ /\ necessarily\ /\ every\ /\ in\ every\ respect} 限定。

\section{逐句练习}\label{ux9010ux53e5ux7ec3ux4e60}

\subsection{S01｜句子拆解}\label{s01ux53e5ux5b50ux62c6ux89e3}

\begin{quote}
Although the court reduced the scope of federal regulation, pollution carried by rivers could still affect communities beyond the states in which it originated.
\end{quote}

\textbf{结构观察}　Although 让步；主句 pollution \ldots{} could affect \ldots；beyond 修饰受影响的范围。

\textbf{最容易丢掉的意义}　把政策边界与实际影响边界分开。

\textbf{动手做}　画出主干，圈出逻辑词，尝试整句翻译。

\subsection{S02｜句子拆解}\label{s02ux53e5ux5b50ux62c6ux89e3}

\begin{quote}
The belief that local laws alone would be sufficient was challenged by the fact that water and the pollutants it carries do not respect administrative boundaries.
\end{quote}

\textbf{结构观察}　主干是 The belief was challenged by the fact。两个 that 从句分别说明 belief 和 fact 的内容；不是把整个名词短语叫作从句。pollutants 后的 it carries 是省略宾语关系代词的定语从句，其中 it 指 water。

\textbf{最容易丢掉的意义}　不要把 that 引出的内容全部当成修饰上一名词的定语从句。

\textbf{动手做}　画出主干，圈出逻辑词，尝试整句翻译。

\subsection{S03｜句子拆解}\label{s03ux53e5ux5b50ux62c6ux89e3}

\begin{quote}
What appears to be a sensible decision for one town may have serious repercussions for people living downstream.
\end{quote}

\textbf{结构观察}　What \ldots{} town 构成主语从句；living downstream 后置修饰 people。

\textbf{最容易丢掉的意义}　主语不是 decision 单词，而是前半个从句所描述的``决定''。

\textbf{动手做}　画出主干，圈出逻辑词，尝试整句翻译。

\subsection{S04｜句子拆解}\label{s04ux53e5ux5b50ux62c6ux89e3}

\begin{quote}
The very measures intended to protect forests from frequent fires may, over decades, increase the fuel available to a later, more destructive blaze.
\end{quote}

\textbf{结构观察}　The very measures 强调；intended to \ldots{} 后置修饰 measures；over decades 是时间限定；available to \ldots{} 修饰 fuel。

\textbf{最容易丢掉的意义}　may 不能译成一定；long-term effects 不否定短期收益。

\textbf{动手做}　画出主干，圈出逻辑词，尝试整句翻译。

\subsection{S05｜句子拆解}\label{s05ux53e5ux5b50ux62c6ux89e3}

\begin{quote}
Only after comparing modern records with evidence preserved in old trees did researchers begin to question whether fewer fires necessarily meant safer forests.
\end{quote}

\textbf{结构观察}　Only after \ldots{} 位于句首引起部分倒装 did researchers begin；还原语序为 researchers began to question \ldots。whether 引出被质疑的命题，preserved in old trees 后置修饰 evidence。

\textbf{最容易丢掉的意义}　句首虽出现助动词 did，但这里是在限定``何时才开始质疑''，不是提出一般疑问句。

\textbf{动手做}　画出主干，圈出逻辑词，尝试整句翻译。

\subsection{S06｜句子拆解}\label{s06ux53e5ux5b50ux62c6ux89e3}

\begin{quote}
Rather than claiming that prescribed burns can eliminate wildfire risk, the report argues that they may be one useful tool when conditions are carefully assessed.
\end{quote}

\textbf{结构观察}　Rather than + V-ing 对比两种主张；that 宾语从句；when \ldots{} 条件/时间语境。

\textbf{最容易丢掉的意义}　不是放任烧火，也不是保证风险归零。

\textbf{动手做}　画出主干，圈出逻辑词，尝试整句翻译。

\subsection{S07｜句子拆解}\label{s07ux53e5ux5b50ux62c6ux89e3}

\begin{quote}
The challenge is not simply to choose between protecting homes and restoring natural processes, but to decide which risks should be addressed immediately and which require longer-term management.
\end{quote}

\textbf{结构观察}　not simply \ldots{} but \ldots{} 是转向结构；which \ldots{} and which \ldots{} 两个平行宾语从句。

\textbf{最容易丢掉的意义}　不能读成只要环境保护就必须牺牲居住安全。

\textbf{动手做}　画出主干，圈出逻辑词，尝试整句翻译。

\subsection{S08｜句子拆解}\label{s08ux53e5ux5b50ux62c6ux89e3}

\begin{quote}
That the number of recorded fires declined is evidence of successful suppression, but it does not by itself establish that the forest became safer in every respect.
\end{quote}

\textbf{结构观察}　That \ldots{} declined 是主语从句；but 转折；it 指前述事实；in every respect 限定推断强度。

\textbf{最容易丢掉的意义}　局部指标变化不等于全面结论。

\textbf{动手做}　画出主干，圈出逻辑词，尝试整句翻译。

\chapter{第四单元｜原创英语一风格阅读}\label{ux7b2cux56dbux5355ux5143ux539fux521bux82f1ux8bedux4e00ux98ceux683cux9605ux8bfb}

以下三篇均为自编教学材料，所述城镇、县与学校是虚构情境。A、B用于同主题巩固，其中B沿用野火管理的论证主题；C转到图书馆服务，检验跨情境迁移。它们不是历年真题，也不是新增实证研究。建议每篇限时8---10分钟完成，暂时不要查词；题目做完后才翻到答案。

\textbf{阅读时请做三件事}：给每段写一个不超过八个字的中文标签；圈出改变段落方向的连接词；选项中任何超出原文范围的词（all / every / never / must）都要重新核查，而不是见到绝对词就机械排除。

\section{Reading A｜Where Flood Protection Begins}\label{reading-awhere-flood-protection-begins}

\emph{城市雨洪治理：政策范围与下游影响}

阅读时先找出谁承担工程费用，再找出为什么作者认为``仅在行政辖区内采取措施''可能不足。

\textbf{1}　For years, a coastal town treated heavy rain as a problem of drainage. When its streets flooded, officials widened pipes and cleared blocked channels. These measures offered visible relief, and residents could see where their taxes had gone. But as storms became more intense, the town discovered a less convenient fact: some of the water that filled its streets had arrived from developments far beyond its borders.

\textbf{2}　The surrounding region had lost patches of grassland and wetland to roads, warehouses and houses. Such development did not necessarily cause flooding at any one site. Yet by reducing the ground's ability to absorb rain, it allowed more runoff to move quickly into streams that passed through the town. The effect of a planning decision was therefore not confined to the jurisdiction in which that decision had been made.

\textbf{3}　Some town councillors called on the regional government to introduce rules limiting construction on remaining wetlands. Builders argued that the proposal was at odds with the demand for affordable housing. They also questioned whether the town had the authority to influence decisions made elsewhere. Both concerns deserved attention. Housing was genuinely needed, and the legal scope of local government could not simply be expanded because a problem was urgent.

\textbf{4}　However, treating the dispute as a choice between building homes and protecting nature missed the larger issue. The costs of a project were not always borne by those who benefited from it. An upstream development could increase flood risks downstream, while the town paying for larger drains had little control over the land-use decisions that made those drains necessary. A regional review would not settle every disagreement, but it could make those connections more visible.

\textbf{5}　The lesson is less about finding a perfect rule than about deciding at what level a problem should be addressed. Pipes remain useful; so do carefully planned housing projects. But where water crosses administrative boundaries, the effects of development should be assessed across those boundaries as well. Ignoring that connection may be the most shortsighted policy of all.

\subsection{Questions}\label{questions}

\Needspace{10\baselineskip}
\textbf{A1.} What did the town discover as storms grew stronger?

A. Most flooding resulted from faulty pipes alone.

B. Some floodwater originated outside the town.

C. Wider streets could prevent damage more effectively.

D. Residents had stopped supporting drainage projects.

\Needspace{10\baselineskip}
\textbf{A2.} The effect of surrounding development is mainly explained by its impact on \_\_\_\_\_.

A. the speed and amount of runoff

B. the level of public housing demand

C. the quality of local tax services

D. the number of regional officials

\Needspace{10\baselineskip}
\textbf{A3.} Why does the author acknowledge the builders' objections?

A. To argue that wetlands provide no useful services.

B. To prove that the town has unlimited regulatory power.

C. To recognize competing needs and legal limits.

D. To recommend abandoning all regional cooperation.

\Needspace{10\baselineskip}
\textbf{A4.} The problem with assessing each development separately is that \_\_\_\_\_.

A. the benefits of new homes cannot be measured

B. the local government cannot collect rainfall data

C. some projects have no direct environmental effects

D. benefits and costs may fall on different communities

\Needspace{10\baselineskip}
\textbf{A5.} Which of the following best expresses the author's view?

A. Problems that cross boundaries require broader assessment.

B. Environmental rules should always replace housing policies.

C. Local drainage improvements are now unnecessary.

D. Any regional review will settle every disagreement.

\textbf{读后复盘}　写出该文的``现象---解释---争议---作者判断''。题目做完后，把你认为最容易误选的一项圈出来，并写出它比原文多说了什么。

\section{Reading B｜When Prevention Has a Cost}\label{reading-bwhen-prevention-has-a-cost}

\emph{森林管理：近期保护与长期风险}

第二篇仍是环境议题，但论证方式不同：先承认某种措施有效，再检视它长期可能造成的副作用。

\textbf{1}　The residents of a mountain county once judged its fire service by a simple measure: the fewer fires, the better. Crews worked quickly to extinguish any blaze near roads, houses or public campsites. For decades, the policy seemed successful. Property losses were limited, and officials could point to a steady decline in the number of recorded fires. Few people questioned a strategy that appeared to offer such clear protection.

\textbf{2}　Then researchers began to compare recent records with evidence preserved in old trees. They found that small fires had once passed through certain parts of the forest far more often than was commonly assumed. Many of these blazes moved slowly and petered out after consuming dry leaves and low vegetation. They did not spare every tree, but they reduced the fuel available to later fires and opened space for new growth. Preventing all such fires was not necessarily the same as making the forest safer.

\textbf{3}　The finding prompted an uncomfortable question. If successful suppression allowed dead branches to accumulate, could past protection have increased some future risks? The answer was not straightforward. Warmer, drier conditions could also allow fires to spread more rapidly, and a single season of extreme weather might overwhelm even well-managed land. Nevertheless, the evidence suggested that judging policy by fire counts alone was shortsighted. Frequency, severity and potential damage were different matters.

\textbf{4}　Officials proposed a limited programme of prescribed burns in carefully chosen locations. Small areas would be treated when wind, humidity and staffing conditions were suitable. Residents welcomed the promise of lower fuel loads but worried about smoke, safety and the possibility of a planned fire escaping control. Those concerns were not merely obstacles to public understanding: they were practical costs that the programme had to address. Officials therefore agreed to publish the conditions under which burns could proceed and to postpone operations when those conditions were not met.

\textbf{5}　The county did not exchange one simple rule for another. It began to measure success by several outcomes rather than one, including the protection of homes, the health of the forest and the risks created by each management choice. That may produce less impressive numbers in the short term. But it offers a more realistic way to balance immediate safety with the ability to withstand future fires.

\subsection{Questions}\label{questions-1}

\Needspace{10\baselineskip}
\textbf{B1.} Why had the county's fire policy seemed successful?

A. It prevented all changes in forest conditions.

B. It eliminated the need for professional fire crews.

C. Recorded fires fell and property losses were limited.

D. It made prescribed burns widely accepted.

\Needspace{10\baselineskip}
\textbf{B2.} The researchers' comparison of old and recent records indicated that \_\_\_\_\_.

A. many small fires once contributed to forest renewal

B. early forests experienced no destructive fires

C. tree-ring evidence was too unreliable to be useful

D. fire activity remained unchanged over time

\Needspace{10\baselineskip}
\textbf{B3.} What does the author suggest about the number of fires?

A. It provides a complete measure of environmental safety.

B. It is irrelevant to evaluating fire policy.

C. It should replace all measures of property damage.

D. It cannot alone show how safe a forest has become.

\Needspace{10\baselineskip}
\textbf{B4.} Residents' concerns about prescribed burns are described as \_\_\_\_\_.

A. proof that the entire plan should be cancelled

B. practical issues to be managed in implementation

C. misunderstandings caused only by poor education

D. evidence that suppression is always preferable

\Needspace{10\baselineskip}
\textbf{B5.} The main lesson of the text is that fire management should \_\_\_\_\_.

A. focus solely on reducing annual fire counts

B. ignore economic costs in favour of natural processes

C. weigh immediate safety and long-term resilience together

D. allow every fire to burn without restriction

\textbf{读后复盘}　写出该文的``现象---解释---争议---作者判断''。题目做完后，把你认为最容易误选的一项圈出来，并写出它比原文多说了什么。

\section{Reading C｜When Fewer Warnings Need a Closer Look}\label{reading-cwhen-fewer-warnings-need-a-closer-look}

\emph{跨情境迁移：图书馆服务与评价指标（自编）}

限时10分钟。先不查词；判断``数字变化''与``它能证明什么''是否被分开。

\textbf{1}　A university library introduced automatic renewal to reduce the number of overdue books. Unless another reader requested an item, its loan period would be extended without any action by the borrower. After one term, overdue notices had fallen sharply. Administrators described the decline as evidence that students were managing their borrowing more responsibly. The new system also saved staff time, allowing them to answer more research enquiries. On these measures, the change looked like an obvious success.

\textbf{2}　Yet a student committee questioned what the figures actually showed. Under the old rules, a book became overdue after a fixed period. Under the new rules, the same book could remain on loan for much longer without generating a notice. The committee did not deny that fewer warnings were welcome. It argued that a change in the definition of a recorded problem should not be confused with a change in the behaviour that had previously produced it.

\textbf{3}　To examine access more directly, the library compiled data on requests, waiting times and unsuccessful searches. Some readers received popular books sooner because the request system brought those books back. Others waited longer for less popular titles that remained with borrowers who no longer needed them. These mixed findings did not establish that automatic renewal had failed. They suggested that the benefits depended partly on how easily readers could discover and request material that was not on the shelves.

\textbf{4}　The committee called on the library to retain automatic renewal while making requests more visible in its catalogue. It also proposed a reminder asking borrowers whether they still needed a book. Some staff worried that extra messages would erode the simplicity students valued. The library therefore tested the reminder with a small group before deciding whether to extend it. A useful tactic in one setting, it reasoned, might create an unnecessary burden in another.

\textbf{5}　The lesson was not that numerical targets should be abandoned. Without records, unequal access might have remained invisible. Rather, institutions needed to ask whether a reported improvement reflected better service, altered recording rules, or both. A policy could curb one inconvenience while creating another; evaluating it required following the experience of users, not simply celebrating the disappearance of an administrative warning.

\subsection{Questions}\label{questions-2}

\Needspace{10\baselineskip}
\textbf{C1.} What does the committee question about the decline in overdue notices?

A. It necessarily indicates more responsible borrowing.

B. It reduced the time available for research enquiries.

C. It followed the introduction of automatic renewal.

D. It was welcomed by some library users.

\Needspace{10\baselineskip}
\textbf{C2.} Why are less popular books mentioned in Paragraph 3?

A. To show that all borrowers prefer popular titles.

B. To establish that loan periods should never be extended.

C. To illustrate differences in the effects on access.

D. To suggest that search records should be discarded.

\Needspace{10\baselineskip}
\textbf{C3.} The proposed reminder was tested with a small group because \_\_\_\_\_.

A. staff had proved that students disliked every message

B. the library wanted to examine a possible trade-off

C. the committee wanted automatic renewal cancelled

D. all books had already returned to the shelves

\Needspace{10\baselineskip}
\textbf{C4.} Which claim is supported by the text?

A. A fall in warnings proves that access has improved for everyone.

B. Changing recording rules cannot affect reported performance.

C. Mixed results make further evaluation unnecessary.

D. Service improvements and changes in recording may occur together.

\Needspace{10\baselineskip}
\textbf{C5.} Which title best captures the author's argument?

A. The Case for Ending All Library Targets

B. Why Borrowers Should Stop Requesting Books

C. When Fewer Warnings Need a Closer Look

D. Automatic Renewal as a Complete Failure

\textbf{迁移任务}　用中文说明：本篇的``逾期通知减少''与野火篇的``火灾次数减少''有什么共同的阅读陷阱？再指出两篇的不同机制，避免机械套模板。

\chapter{第五单元｜练习：让词真正进入理解}\label{ux7b2cux4e94ux5355ux5143ux7ec3ux4e60ux8ba9ux8bcdux771fux6b63ux8fdbux5165ux7406ux89e3}

\section{A. 语境选词填空（V01---V16）}\label{a.-ux8bedux5883ux9009ux8bcdux586bux7a7av01v16}

从词库中选取合适词语，必要时变化词形。每个表达只用一次，短语须完整填入。先做题再对答案。

\textbf{词库}　compile / jurisdiction / erode / runoff / scope / call on / shortfall / authority / peter out / shortsighted / prescribed / repercussion / curb / at odds with / suppression / transcend

\textbf{V01.}　The committee will \_\_\_\_\_\_ information from several agencies before reaching a decision.

\textbf{V02.}　The river runs through three separate \_\_\_\_\_\_, making joint management difficult.

\textbf{V03.}　Over the past year, a lack of investment has gradually \_\_\_\_\_\_ confidence in the service.

\textbf{V04.}　Heavy rain carried polluted \_\_\_\_\_\_ from farms into the lake.

\textbf{V05.}　The study lies outside the \_\_\_\_\_\_ of the current review.

\textbf{V06.}　Yesterday, the editorial \_\_\_\_\_\_ officials to consider the interests of nearby residents.

\textbf{V07.}　The budget \_\_\_\_\_\_ led to a delay in repairing damaged bridges.

\textbf{V08.}　The new law gave the agency the \_\_\_\_\_\_ to inspect unsafe buildings.

\textbf{V09.}　The small fire \_\_\_\_\_\_ when it reached a patch of wet ground.

\textbf{V10.}　The plan was criticized as \_\_\_\_\_\_ because it ignored maintenance costs.

\textbf{V11.}　The fire service carried out a \_\_\_\_\_\_ burn under suitable weather conditions.

\textbf{V12.}　The decision could have unexpected \_\_\_\_\_\_ for towns downstream.

\textbf{V13.}　New measures were introduced to \_\_\_\_\_\_ the spread of the disease.

\textbf{V14.}　The two departments appeared to be \_\_\_\_\_\_ one another.

\textbf{V15.}　Several years of fire \_\_\_\_\_\_ increased the amount of dry wood on the forest floor.

\textbf{V16.}　An agreement may help \_\_\_\_\_\_ the political divisions that hinder cooperation.

\section{B. 一词多义与程度选择（M01---M06）}\label{b.-ux4e00ux8bcdux591aux4e49ux4e0eux7a0bux5ea6ux9009ux62e9m01m06}

判断下列句中加粗词或整个表达最接近哪一个英文解释。中文提示在附录，先不要看。

\Needspace{10\baselineskip}
\textbf{M01.}　The law limits the agency's authority, not the river's flow.

A. legal power or permission

B. an expert respected for knowledge

\Needspace{10\baselineskip}
\textbf{M02.}　The court changed the scope of regulation, but the ecological consequences may unfold later.

A. all environmental rules have been removed

B. the range of regulation has changed

\Needspace{10\baselineskip}
\textbf{M03.}　A prescribed burn is not the same as allowing any wildfire to spread unchecked.

A. a carefully planned fire

B. any fire allowed to spread freely

\Needspace{10\baselineskip}
\textbf{M04.}　A decline in the number of fires does not establish that future losses will decline.

A. fewer fires prove all risks are lower

B. fewer fires do not prove every risk is lower

\Needspace{10\baselineskip}
\textbf{M05.}　The findings were a testament to the benefits of careful monitoring.

A. evidence or a clear sign of something

B. a legal document about inheritance

\Needspace{10\baselineskip}
\textbf{M06.}　The policy is a double-edged sword: it reduces one risk while creating another.

A. a policy free of trade-offs

B. a policy with both benefits and disadvantages

\section{C. 推理与限定（I01---I04）}\label{c.-ux63a8ux7406ux4e0eux9650ux5b9ai01i04}

不是判断现实中事情是否可能，而是判断哪个选项得到了\textbf{给定句子}的支持。

\Needspace{10\baselineskip}
\textbf{I01.}　The wetland review may affect several counties, although only one county requested it.

\emph{The sentence allows which inference?}

A. The review must be limited to the requesting county.

B. The review has already changed every local law.

C. The possible effects extend beyond the requesting county.

D. No county has a legal interest in the review.

\Needspace{10\baselineskip}
\textbf{I02.}　The measure reduced flooding in the town centre but increased runoff into a nearby village.

\emph{What is the best conclusion?}

A. The project had different effects in different places.

B. The project eliminated flooding across the region.

C. The village received no additional runoff.

D. The project was necessarily unlawful.

\Needspace{10\baselineskip}
\textbf{I03.}　The study suggests that a policy can succeed by one measure while failing by another.

\emph{The main idea is that \_\_\_\_\_.}

A. one indicator always determines success

B. policies should never be evaluated

C. all results are equally important

D. policy outcomes may depend on how success is measured

\Needspace{10\baselineskip}
\textbf{I04.}　Prescribed burns may help reduce fuel, but they must be postponed when weather conditions are unsuitable.

\emph{Which statement is consistent with this?}

A. Fire managers can burn whenever they choose.

B. The method has potential benefits but requires limits.

C. Prescribed burns never produce smoke.

D. Poor weather makes all forest management useless.

\section{D. 英译汉专项（T01---T06）}\label{d.-ux82f1ux8bd1ux6c49ux4e13ux9879t01t06}

保留否定、可能性、指代和范围。不要只追求中文通顺，而丢掉作者刻意保留的限定。

\textbf{T01.}　A reduction in federal authority does not automatically remove every form of protection available under state law.

译文：\_\_\_\_\_\_\_\_\_\_\_\_\_\_\_\_\_\_\_\_\_\_\_\_\_\_\_\_\_\_\_\_\_\_\_\_\_\_\_\_

\textbf{T02.}　The costs of pollution may be felt far from the place where decisions about land use were originally made.

译文：\_\_\_\_\_\_\_\_\_\_\_\_\_\_\_\_\_\_\_\_\_\_\_\_\_\_\_\_\_\_\_\_\_\_\_\_\_\_\_\_

\textbf{T03.}　It would be shortsighted to measure a policy only by its immediate benefits when the effects of that policy may last for decades.

译文：\_\_\_\_\_\_\_\_\_\_\_\_\_\_\_\_\_\_\_\_\_\_\_\_\_\_\_\_\_\_\_\_\_\_\_\_\_\_\_\_

\textbf{T04.}　A small fire that clears dry vegetation may reduce the intensity of later fires without preventing them entirely.

译文：\_\_\_\_\_\_\_\_\_\_\_\_\_\_\_\_\_\_\_\_\_\_\_\_\_\_\_\_\_\_\_\_\_\_\_\_\_\_\_\_

\textbf{T05.}　The fact that a programme has lowered the number of reported fires is not sufficient evidence that it has improved safety in every respect.

译文：\_\_\_\_\_\_\_\_\_\_\_\_\_\_\_\_\_\_\_\_\_\_\_\_\_\_\_\_\_\_\_\_\_\_\_\_\_\_\_\_

\textbf{T06.}　Even where a planned burn is considered beneficial, the timing of the operation must depend on local weather and available personnel.

译文：\_\_\_\_\_\_\_\_\_\_\_\_\_\_\_\_\_\_\_\_\_\_\_\_\_\_\_\_\_\_\_\_\_\_\_\_\_\_\_\_

\section{E. 句子输出（P01---P05）}\label{e.-ux53e5ux5b50ux8f93ux51fap01p05}

用已经学过的词或短语表达下面的中文。重点检查动词和宾语是否能够自然搭配。参考译法放在答案部分。

\textbf{P01.}　用 scope of 表达：报告没有讨论所有环境问题。

英文：\_\_\_\_\_\_\_\_\_\_\_\_\_\_\_\_\_\_\_\_\_\_\_\_\_\_\_\_\_\_\_\_\_\_\_\_\_\_\_\_

\textbf{P02.}　用 call on sb. to do 表达：研究者呼吁地方政府分享水质数据。

英文：\_\_\_\_\_\_\_\_\_\_\_\_\_\_\_\_\_\_\_\_\_\_\_\_\_\_\_\_\_\_\_\_\_\_\_\_\_\_\_\_

\textbf{P03.}　用 erode 与 may 表达：长期资金不足可能逐渐削弱公众信心。

英文：\_\_\_\_\_\_\_\_\_\_\_\_\_\_\_\_\_\_\_\_\_\_\_\_\_\_\_\_\_\_\_\_\_\_\_\_\_\_\_\_

\textbf{P04.}　用 at odds with 表达：该决定似乎与现有证据不一致。

英文：\_\_\_\_\_\_\_\_\_\_\_\_\_\_\_\_\_\_\_\_\_\_\_\_\_\_\_\_\_\_\_\_\_\_\_\_\_\_\_\_

\textbf{P05.}　用 rather than 表达：应比较多项结果，而不是只计算火灾次数。

英文：\_\_\_\_\_\_\_\_\_\_\_\_\_\_\_\_\_\_\_\_\_\_\_\_\_\_\_\_\_\_\_\_\_\_\_\_\_\_\_\_

\chapter{第六单元｜做完这册，你应该能够}\label{ux7b2cux516dux5355ux5143ux505aux5b8cux8fd9ux518cux4f60ux5e94ux8be5ux80fdux591f}

不看释义，尝试独立回答下面六个问题。这不是给自己打分，而是核对本册最值得带走的语言能力。

\begin{enumerate}
\def\labelenumi{\arabic{enumi}.}
\tightlist
\item
  \texttt{authority\ /\ jurisdiction\ /\ scope} 各自强调什么？能否各造一条短句？
\item
  2024年那篇文章中，湿地功能与法律裁决之间有什么关系？为什么作者要谈流域？
\item
  \texttt{curb\ /\ suppress\ /\ eliminate} 的力度有什么不同？
\item
  为什么2026文章里的 fire shortfall 不能读成``全世界的野火损失都下降''？
\item
  读到 \texttt{may\ /\ not\ necessarily\ /\ all\ /\ only}，怎样防止选项改变原文的范围？
\item
  怎样区分``某措施的目标''``观察到的结果''和``作者的评价''？
\end{enumerate}

把回答不清楚的点做一个小标记，下次阅读时再练。你不需要把所有句子背下来，但应能在另一篇文章中识别同一种表达与推理。

\chapter{答案与解析}\label{ux7b54ux6848ux4e0eux89e3ux6790}

\section{一、开始前的诊断}\label{ux4e00ux5f00ux59cbux524dux7684ux8bcaux65ad}

\textbf{D01 \(\rightarrow\) B}　scope 指覆盖范围，不是湿地实际大小。 A把制度范围换成实际面积；C是天气量；D是政府数量，后三者均未在句中改变。

\textbf{D02 \(\rightarrow\) C}　at odds with 直接给出不一致的关系。 A颠倒关系；B把已表明的冲突弱化成不确定；D是重复而非不一致。

\textbf{D03 \(\rightarrow\) A}　curb 与 eliminate 不等价；does not necessarily 是否定必然性。 B把部分作用抹成无作用；C把未必清除全部改成已经全部消失；D反转作用方向并增加总是。

\textbf{D04 \(\rightarrow\) D}　may 是可能性；短期目的不保证长期后果。 A否认了可能的短期收益；B把 may 升级成必然；C引入未提到的立法难度。

\section{二、2024年真题 Text 4（36---40）}\label{ux4e8c2024ux5e74ux771fux9898-text-43640}

\textbf{36 A}　首段通过生境、蓄洪、过滤污染物和防护风暴潮等不同功能说明湿地环境价值。不要把``多种生态功能''窄化为单一商业收益。

\textbf{37 B}　法律裁决收缩的是EPA监管适用范围。对照原文的 \texttt{authority} 和 \texttt{regulatory\ scope}。A把裁决方向反转为加强污染控制；C凭空增加居民冲突终结；D把受益于裁决的矿业经营者写成可能反对者。

\textbf{38 A}　作者对切萨皮克湾的未来感到担忧。抓住 \texttt{shortsighted} 和关于其他辖区是否同样保护湿地的质疑。并非``困惑''或``已经放心''。

\textbf{39 D}　跨州水流导致上游行为可能影响下游，EPA长期参与对协调保护有重要作用。正确项概括协调功能，而非声称它``已经恢复所有辖区之间的平衡''。

\textbf{40 B}　结尾是敦促州立法者认真对待湿地保护。\texttt{call\ on} 提示行动发起者，\texttt{lawmakers} 指行动对象；A把行动重点移到土地所有者的影响；C、D抓住末段保护区和生物的细节，却未概括审视湿地保护的政策请求。

以上字母来自公开的非官方参考解析，作学习核对用。请以你打开的试卷上的36---40题为题面。

\section{三、2026年真题 Text 4（36---40）}\label{ux4e092026ux5e74ux771fux9898-text-43640}

\textbf{36 D}　研究者惊讶于殖民者到来之前野火频繁发生的记录。1970年代的职责只是引出旧观念的背景，不是``惊讶之处''。

\textbf{37 C}　第2段集中说明低烈度火灾的生态收益，例如清理枯木、降低竞争与病害传播风险。不能把列举中的``病树''误当整段唯一主题。

\textbf{38 B}　研究点位近期火灾发生频率低于依据历史节律作出的比较预期。这里是\textbf{频率}，不是强度、死亡人数或全球火灾总损失。

\textbf{39 D}　长期压制小火和气候变化等因素可能使极端火更难扑灭；推断保留了 \texttt{may}，没有把现代扑火说成完全无用。

\textbf{40 C}　两位研究者共同支持在适当条件下开展计划烧除，也承认气候变化增加实施难度。共同立场是开展计划烧除的必要性。A把实施困难换成预测困难；B把实际采用换成紧迫评估；D增加了加倍火灾侦测的主张。

以上亦为公开非官方参考答案。

\section{四、原创阅读｜逐题解析}\label{ux56dbux539fux521bux9605ux8bfbux9010ux9898ux89e3ux6790}

\subsection{Reading A｜答案与选项比较}\label{reading-aux7b54ux6848ux4e0eux9009ux9879ux6bd4ux8f83}

\textbf{A1 \(\rightarrow\) B}　第一段的转向不否定已有排水措施，而是补充了``外部来水''。A把原因缩成单一管道问题；C无据；D没有居民反对的事实。

\textbf{A2 \(\rightarrow\) A}　第二段有完整的因果链：植被减少\(\rightarrow\)地表吸水下降\(\rightarrow\)径流更快更多\(\rightarrow\)下游承压。正确项概括了中间机制，其余项与解释无关。 逐项看：B用住房需求替代水文机制；C的税收服务未被讨论；D引入官员人数。

\textbf{A3 \(\rightarrow\) C}　第三段并未把一方妖魔化，而是承认住房需求与法律权限问题。C准确兼顾两个限定；A、B、D是过度化。 A否定湿地价值，B声称无限权限，均与文意相反；D要求放弃地区合作，而文章正解释为何需要跨区评估。

\textbf{A4 \(\rightarrow\) D}　第四段强调外部性：上游开发，下游受损，承担修管费用者未必能影响源头决策。D概括这一结构；A、B、C均不属于段落主旨。 A和B分别增加收益无法测量及不能收集雨量的说法；C谈有无直接影响，遗漏成本转移这一中心关系。

\textbf{A5 \(\rightarrow\) A}　末段明确保留``排水工程和住房仍有价值''，主张扩展影响评估的尺度；B/C/D分别是绝对化、否定原有措施与夸大作用。

\subsection{Reading B｜答案与选项比较}\label{reading-bux7b54ux6848ux4e0eux9009ux9879ux6bd4ux8f83}

\textbf{B1 \(\rightarrow\) C}　第一段明确给出两项衡量结果：火灾记录下降、财产损失有限。它们说明过去政策具有现实收益，但不证明长期风险也最小。 A把财产损失有限扩大为森林毫无变化；B与专业队伍积极扑火冲突；D把后文的新政策移到过去。

\textbf{B2 \(\rightarrow\) A}　第二段列出低烈度火的生态作用，A是谨慎概括；B/C/D均与文章事实冲突。 B把许多小火的益处夸成从无破坏性火灾；C否定文章实际使用的证据；D否定历史与近期记录的差异。

\textbf{B3 \(\rightarrow\) D}　第三段明确把频率、严重程度与损失区分开，D准确表达``单一指标不足''；B又把``不充分''扭成``完全没用''。 A把单一指标说成充分；C要求取代损失指标，文中没有这种主张。

\textbf{B4 \(\rightarrow\) B}　第四段先列风险，再说明具体的条件限制和延期措施。作者既未否认风险，也未主张因此彻底放弃计划。 A把执行风险升级成必须取消；C将实际风险全归为教育不足；D在有限情境外添加始终更好。

\textbf{B5 \(\rightarrow\) C}　最后一段从单一指标转向多维评价，强调眼前安全、森林健康和未来风险；C保留了完整限定。 A仍只算次数；B忽略成本，而文章要求权衡；D从有计划干预变成所有火不受限制地燃烧。

\subsection{Reading C｜答案与证据}\label{reading-cux7b54ux6848ux4e0eux8bc1ux636e}

\textbf{C1 \(\rightarrow\) A}　P1---P2。A是委员会质疑的推断，而非作者认可的事实：通知减少也可能来自记账规则变化。B与节省工作人员时间相反；C是双方接受的先后事实；D并非争议焦点，委员会明确不否认减少警告的好处。

\textbf{C2 \(\rightarrow\) C}　P3对照热门书与冷门书的等待时间。C涵盖两类结果；A把一个例子扩大成所有人的偏好；B把某些人的等待延长升级为一律禁止；D与收集数据来发现差异的做法相反。

\textbf{C3 \(\rightarrow\) B}　P4同时提出促进归还的设想与额外信息负担的担忧。B保留``需要考察''的不确定性；A把担忧写成证明并扩大到所有消息；C与保留自动续借的建议相反；D没有依据。

\textbf{C4 \(\rightarrow\) D}　P5明确要求区分服务、记录规则及二者共同作用。D不把二者设为互斥；A从单一指标推出人人获益；B否认了P2的机制；C把继续调查的理由变成停止调查的理由。

\textbf{C5 \(\rightarrow\) C}　C串起指标变化、机制核查与使用者体验。A被末段开头明确排除；B颠倒了让预约更方便的建议；D将不均衡结果夸大成彻底失败，抹去了节省时间和部分等待缩短的收益。

\textbf{迁移任务参考}　共同点是单一数字不足以推出总体改善。不同点在于，图书馆篇改变了触发通知的规则，还涉及借书可获得性；野火篇则涉及长期扑火、燃料积累与火灾后果。不能把野火次数变化也解释为记录规则改变。

\section{五、选词填空}\label{ux4e94ux9009ux8bcdux586bux7a7a}

\textbf{V01}　compile。先有分散材料，后汇集整理。

\textbf{V02}　jurisdictions。注意复数。

\textbf{V03}　eroded。has 后用过去分词构成现在完成时；gradually 与逐渐削弱的含义呼应。

\textbf{V04}　runoff。可作不可数名词，强调径流。

\textbf{V05}　scope。固定结构 outside the scope of。

\textbf{V06}　called on。call on sb. to do。

\textbf{V07}　shortfall。shortfall 是有比较基准的不足。

\textbf{V08}　authority。法律赋权，不是专家权威。

\textbf{V09}　petered out。固定短语 peter out。

\textbf{V10}　shortsighted。作者评价语，注意原因。

\textbf{V11}　prescribed。固定专业搭配 prescribed burn。

\textbf{V12}　repercussions。常用复数。

\textbf{V13}　curb。强调抑制不是完全根除。

\textbf{V14}　at odds with。at odds with 固定短语。

\textbf{V15}　suppression。不等于所有形式的森林管理。

\textbf{V16}　transcend。搭配政治分歧等抽象边界。

\section{六、词义判断}\label{ux516dux8bcdux4e49ux5224ux65ad}

\textbf{M01 \(\rightarrow\) A}　此处指机构的法律权限；B把权限义换成了``权威专家''义，并非句中所指。

\textbf{M02 \(\rightarrow\) B}　scope 的变化是适用范围变化；A把范围调整扩大成取消全部规则。句子也没有指定监管属于联邦层级。

\textbf{M03 \(\rightarrow\) A}　计划与条件是关键区别。 B删去了方案和条件，把受约束的活动改成任意蔓延。

\textbf{M04 \(\rightarrow\) B}　频率不能直接推出损失。 A把一个指标的下降扩大成全部风险降低。

\textbf{M05 \(\rightarrow\) A}　testament 表示见证/体现，不是遗嘱。 B属于另一个词义，与 findings 和 monitoring 的证据语境不合。

\textbf{M06 \(\rightarrow\) B}　both sides 同时存在。 A宣称没有权衡，与一边降低风险、一边增加风险直接矛盾。

\section{七、限定与推断}\label{ux4e03ux9650ux5b9aux4e0eux63a8ux65ad}

\textbf{I01 \(\rightarrow\) C}　范围词 only 限定提出请求者，不限定受影响者；may 保持可能性。 A错把 only 的范围转到影响上；B把可能写成既成事实并加上每项法律；D无依据地否认法律利益。

\textbf{I02 \(\rightarrow\) A}　地点差异被明确提到，不能升级为整体消除或违法。 B抹去村庄受到的不利影响；C直接否认新增径流；D增加法律定性。

\textbf{I03 \(\rightarrow\) D}　强调评价标准影响判断，不主张一切标准等权。 A忽视不同标准；B否认评估必要性；C从标准不同推出所有结果同等重要，证据不足。

\textbf{I04 \(\rightarrow\) B}　but 引入执行前提，may 不代表保证。 A取消天气限制；C承诺从不冒烟，原句没说；D把一种操作暂缓扩大为所有管理无用。

\section{八、英译汉参考译文}\label{ux516bux82f1ux8bd1ux6c49ux53c2ux8003ux8bd1ux6587}

\textbf{T01}　联邦权限的缩减，并不自动取消州法下所有可用的保护措施。

\emph{提示}：范围否定：not automatically 和 every form 都不能漏译。

\textbf{T02}　污染的代价可能在距离土地利用决策最初作出地点很远的地方显现。

\emph{提示}：where 引导地点关系；felt 为被动意义。

\textbf{T03}　如果一项政策的影响可能持续数十年，却只用眼前收益评价它，这将是目光短浅的。

\emph{提示}：when 可按具体语义译作``在\ldots\ldots 情况下/如果''；后半句承载论证条件。

\textbf{T04}　一场清除干燥植被的小火，可能降低后续火灾的强度，却未必能彻底阻止它们发生。

\emph{提示}：without preventing them entirely 并非``不会阻止任何火灾''。

\textbf{T05}　一项计划减少了报告中的火灾次数，并不足以证明它已经在所有方面提高了安全程度。

\emph{提示}：the fact that \ldots{} 作主语；not sufficient evidence that \ldots{} 注意推断边界。

\textbf{T06}　即使计划烧除被认为有益，实施时间仍须取决于当地天气和可用人员。

\emph{提示}：Even where 意为``即便在\ldots\ldots 情况下''；must depend on 体现限制。

\section{九、汉译英参考写法}\label{ux4e5dux6c49ux8bd1ux82f1ux53c2ux8003ux5199ux6cd5}

\textbf{P01}　Not every environmental issue falls within the scope of the report.

\emph{提示}：可采用更自然版本：The report does not address every environmental issue; some fall outside its scope.

\textbf{P02}　The researchers called on local authorities to share water-quality data.

\emph{提示}：动作发出者是研究者；执行者是地方政府。

\textbf{P03}　Persistent underfunding may gradually erode public confidence.

\emph{提示}：保留``可能''，不能升级为必然。

\textbf{P04}　The decision appears to be at odds with the available evidence.

\emph{提示}：appears to be 表示从表面证据作出的判断。

\textbf{P05}　We should compare several outcomes rather than count fires alone.

\emph{提示}：rather than 后面保持两个并列动词结构。

\section{十、精读停顿与自检参考}\label{ux5341ux7cbeux8bfbux505cux987fux4e0eux81eaux68c0ux53c2ux8003}

\textbf{停顿（1）}　前提是本州有法律；遗漏条件是域外活动同样可能影响水质；不能直接得出的结论是跨州保护已经充分。作者没有否定州法，而是否定由局部保障推出整体安全。

\textbf{停顿（2）}　上游土地利用影响地表径流，径流携带污染物进入水系，下游社区承担影响；成本与决策地点分离，所以行政辖区内的单独管理可能不足。

\textbf{停顿（3）}　20÷100=20\%，估计差额为80次。不能据此推出经济损失或死亡人数下降80\%，因为题目没有给出这些指标，参照值也属于假设估计。

\textbf{停顿（4）}　可答：现代灭火降低了研究点位的火灾次数；长期压制部分小火可能增加燃料积累，且气候变化也影响风险；计划烧除必须考虑适合的天气与实施条件。

\textbf{第一单元``把词放回文章''}　问题一：法定权限是制度变量，污染下降须用环境观测支持。问题二：违法需有适用法律及违反行为的证据，不一致本身不够。问题三：裁决改变监管边界，而径流说明这种边界变化为何会有域外后果。问题四：批评的是局部自足的保护观；跨州流域提供论据。

\textbf{第六单元自检}　（1）authority偏重法定权限，jurisdiction兼指管辖权及辖区，scope指适用范围，三者并非严格互斥。（2）生态价值解释保护为何重要；流域联系解释为何不能只看本州。（3）curb是遏制，suppress在本篇是压制或扑灭，eliminate是消除，具体结果仍受上下文限制。（4）缺口是研究点位相对参照的发生次数不足，没有直接测量全球损失。（5）逐个检查可能性、否定作用域、数量及排他范围。（6）目的写``旨在''，结果写实际观察，评价必须说明采用的标准。

\chapter{英中双语对照附录}\label{ux82f1ux4e2dux53ccux8bedux5bf9ux7167ux9644ux5f55}

这里集中提供本册原创英文例句、原创阅读全文、题干和选项、练习句子的中译。\textbf{安排在答案之后，是为了让第一次做题时仍然能够独立依靠英文完成理解。} 中文力求忠实保留程度、时态、否定与条件。

\section{附录A｜重点词例句}\label{ux9644ux5f55aux91cdux70b9ux8bcdux4f8bux53e5}

\Needspace{5\baselineskip}
\textbf{01 filter}

\begin{samepage}
\textbf{01-1}　Wetland vegetation can filter some pollutants before they reach open water.\\
\zhline{湿地植被能在部分污染物进入开阔水域之前将其过滤。}
\par\end{samepage}

\begin{samepage}
\textbf{01-2}　The committee filtered out unreliable reports before reviewing the evidence.\\
\zhline{委员会在审查证据之前筛除了不可靠的报告。}
\par\end{samepage}

\Needspace{5\baselineskip}
\textbf{02 authority}

\begin{samepage}
\textbf{02-1}　The agency has the authority to inspect factories that discharge waste into the river.\\
\zhline{该机构有权检查向河流排放废物的工厂。}
\par\end{samepage}

\begin{samepage}
\textbf{02-2}　Several local authorities expressed concern about the proposed changes.\\
\zhline{若干地方当局对拟议的变化表示担忧。}
\par\end{samepage}

\Needspace{5\baselineskip}
\textbf{03 scope}

\begin{samepage}
\textbf{03-1}　The decision narrowed the scope of the regulation without abolishing it.\\
\zhline{这项决定缩小了法规的适用范围，但没有废除法规。}
\par\end{samepage}

\begin{samepage}
\textbf{03-2}　The report is valuable, although several issues fall beyond its scope.\\
\zhline{这份报告很有价值，不过若干问题超出了它的讨论范围。}
\par\end{samepage}

\Needspace{5\baselineskip}
\textbf{04 at odds with}

\begin{samepage}
\textbf{04-1}　The policy is at odds with the evidence gathered by independent researchers.\\
\zhline{这项政策与独立研究者收集到的证据不一致。}
\par\end{samepage}

\begin{samepage}
\textbf{04-2}　Her public statement seemed at odds with the company's earlier promises.\\
\zhline{她的公开声明似乎与公司先前的承诺相矛盾。}
\par\end{samepage}

\Needspace{5\baselineskip}
\textbf{05 repercussion}

\begin{samepage}
\textbf{05-1}　A change in upstream farming may have repercussions for communities downstream.\\
\zhline{上游耕作方式的变化可能给下游社区带来连带影响。}
\par\end{samepage}

\begin{samepage}
\textbf{05-2}　The decision had unintended repercussions for small businesses.\\
\zhline{这项决定给小企业带来了意料之外的连锁影响。}
\par\end{samepage}

\Needspace{5\baselineskip}
\textbf{06 jurisdiction}

\begin{samepage}
\textbf{06-1}　The river crosses several jurisdictions, each with different environmental rules.\\
\zhline{这条河流经若干管辖区域，而各地的环保规则并不相同。}
\par\end{samepage}

\begin{samepage}
\textbf{06-2}　The dispute concerns whether the case falls within federal jurisdiction.\\
\zhline{争议在于这一案件是否属于联邦管辖范围。}
\par\end{samepage}

\Needspace{5\baselineskip}
\textbf{07 shortsighted}

\begin{samepage}
\textbf{07-1}　It would be shortsighted to ignore maintenance costs when choosing a cheaper system.\\
\zhline{选择较便宜的系统时若忽视维护成本，就显得缺乏长远考虑。}
\par\end{samepage}

\begin{samepage}
\textbf{07-2}　A shortsighted decision may bring savings now but greater costs later.\\
\zhline{目光短浅的决定可能带来眼前节省，却造成更大的后续成本。}
\par\end{samepage}

\Needspace{5\baselineskip}
\textbf{08 transcend}

\begin{samepage}
\textbf{08-1}　Effective river management must transcend administrative boundaries.\\
\zhline{有效的河流治理必须超越行政区划的限制。}
\par\end{samepage}

\begin{samepage}
\textbf{08-2}　The exhibition aims to transcend differences between generations.\\
\zhline{这场展览旨在跨越不同世代之间的隔阂。}
\par\end{samepage}

\Needspace{5\baselineskip}
\textbf{09 runoff}

\begin{samepage}
\textbf{09-1}　Heavy rain generated runoff that carried soil from fields into nearby streams.\\
\zhline{暴雨形成的地表径流将农田中的泥土带入附近的小溪。}
\par\end{samepage}

\begin{samepage}
\textbf{09-2}　Vegetated buffer strips can filter pollutants from runoff before it reaches the lake.\\
\zhline{植被缓冲带能在径流进入湖泊之前滤除其中的污染物。}
\par\end{samepage}

\Needspace{5\baselineskip}
\textbf{10 call on}

\begin{samepage}
\textbf{10-1}　The editorial called on lawmakers to review existing wetland protections.\\
\zhline{这篇社论呼吁立法者审查现行湿地保护措施。}
\par\end{samepage}

\begin{samepage}
\textbf{10-2}　The researchers called on schools to reconsider how success is measured.\\
\zhline{研究者呼吁学校重新考虑成功应如何衡量。}
\par\end{samepage}

\Needspace{5\baselineskip}
\textbf{11 stakes}

\begin{samepage}
\textbf{11-1}　The stakes are high for villages that depend on the river for drinking water.\\
\zhline{对于依赖这条河获得饮用水的村庄而言，利害关系重大。}
\par\end{samepage}

\begin{samepage}
\textbf{11-2}　What is at stake is not just a single species but the stability of the wider ecosystem.\\
\zhline{面临风险的不仅是某一个物种，还有更广泛生态系统的稳定性。}
\par\end{samepage}

\Needspace{5\baselineskip}
\textbf{12 mandate}

\begin{samepage}
\textbf{12-1}　The forest service received a mandate to protect nearby communities from fire.\\
\zhline{森林管理机构被赋予保护附近社区免受火灾威胁的职责。}
\par\end{samepage}

\begin{samepage}
\textbf{12-2}　A new law may mandate regular safety inspections for public buildings.\\
\zhline{新法律可能规定公共建筑必须定期接受安全检查。}
\par\end{samepage}

\Needspace{5\baselineskip}
\textbf{13 curb}

\begin{samepage}
\textbf{13-1}　Small fires can help curb the spread of some forest diseases.\\
\zhline{小规模火灾有助于遏制某些森林病害的传播。}
\par\end{samepage}

\begin{samepage}
\textbf{13-2}　The city introduced measures to curb excessive water use during the drought.\\
\zhline{这座城市采取措施限制干旱期间过量用水。}
\par\end{samepage}

\Needspace{5\baselineskip}
\textbf{14 peter out}

\begin{samepage}
\textbf{14-1}　The fire petered out when it reached an area with little dry vegetation.\\
\zhline{火势到达缺少干燥植被的地带后逐渐熄灭。}
\par\end{samepage}

\begin{samepage}
\textbf{14-2}　Public interest in the campaign petered out after several months.\\
\zhline{几个月后，公众对这场活动的兴趣渐渐消退。}
\par\end{samepage}

\Needspace{5\baselineskip}
\textbf{15 compile}

\begin{samepage}
\textbf{15-1}　The scientists compiled fire records from several research sites.\\
\zhline{科学家汇集了多个研究地点的火灾记录。}
\par\end{samepage}

\begin{samepage}
\textbf{15-2}　The museum compiled visitors' comments into an annual report.\\
\zhline{博物馆将参观者的意见整理成一份年度报告。}
\par\end{samepage}

\Needspace{5\baselineskip}
\textbf{16 shortfall}

\begin{samepage}
\textbf{16-1}　The records revealed a shortfall in small fires compared with the number expected at the site.\\
\zhline{记录显示，该地点发生的小规模火灾次数少于预期。}
\par\end{samepage}

\begin{samepage}
\textbf{16-2}　A funding shortfall forced the town to delay the restoration project.\\
\zhline{经费缺口迫使小镇推迟修复项目。}
\par\end{samepage}

\Needspace{5\baselineskip}
\textbf{17 suppression}

\begin{samepage}
\textbf{17-1}　Decades of fire suppression allowed dead wood to accumulate.\\
\zhline{数十年的火灾压制使枯木不断积累。}
\par\end{samepage}

\begin{samepage}
\textbf{17-2}　The report examines whether strict suppression remains the best policy.\\
\zhline{报告考察了严格扑火是否仍是最佳政策。}
\par\end{samepage}

\Needspace{5\baselineskip}
\textbf{18 testament}

\begin{samepage}
\textbf{18-1}　The survival of ancient trees is a testament to the forest's resilience.\\
\zhline{古树得以存活，体现了森林的恢复力。}
\par\end{samepage}

\begin{samepage}
\textbf{18-2}　The new library is a testament to years of community cooperation.\\
\zhline{新图书馆是社区多年合作成果的见证。}
\par\end{samepage}

\Needspace{5\baselineskip}
\textbf{19 erode}

\begin{samepage}
\textbf{19-1}　A growing fuel load can erode firefighters' ability to control severe fires.\\
\zhline{不断增加的可燃物可能逐渐削弱消防人员控制严重火灾的能力。}
\par\end{samepage}

\begin{samepage}
\textbf{19-2}　Repeated delays eroded public confidence in the transport plan.\\
\zhline{接连拖延逐渐削弱了公众对交通计划的信心。}
\par\end{samepage}

\Needspace{5\baselineskip}
\textbf{20 ignite}

\begin{samepage}
\textbf{20-1}　Strong winds can help a small spark ignite dry grass.\\
\zhline{大风可能使小火星点燃干草。}
\par\end{samepage}

\begin{samepage}
\textbf{20-2}　The proposal ignited a debate over who should pay for flood protection.\\
\zhline{这项提议引发了关于谁该为防洪买单的争论。}
\par\end{samepage}

\Needspace{5\baselineskip}
\textbf{21 prescribed}

\begin{samepage}
\textbf{21-1}　Prescribed burns are conducted under carefully selected conditions.\\
\zhline{计划烧除是在经过审慎选择的条件下进行的。}
\par\end{samepage}

\begin{samepage}
\textbf{21-2}　The prescribed procedure must be followed when dangerous materials are handled.\\
\zhline{处理危险材料时必须遵守规定程序。}
\par\end{samepage}

\Needspace{5\baselineskip}
\textbf{22 tactic}

\begin{samepage}
\textbf{22-1}　Careful thinning is one tactic for reducing the amount of fuel near homes.\\
\zhline{审慎疏伐是减少住宅附近可燃物的一种办法。}
\par\end{samepage}

\begin{samepage}
\textbf{22-2}　The campaign adopted a different tactic after its early efforts failed.\\
\zhline{前期尝试失败后，这场活动采用了另一种策略。}
\par\end{samepage}

\Needspace{5\baselineskip}
\textbf{23 double-edged sword}

\begin{samepage}
\textbf{23-1}　Greater public access can be a double-edged sword for nature reserves.\\
\zhline{扩大公众进入自然保护区的机会，可能是一把双刃剑。}
\par\end{samepage}

\begin{samepage}
\textbf{23-2}　Digital monitoring is a double-edged sword when privacy safeguards are weak.\\
\zhline{当隐私保护不足时，数字监测就成为一把双刃剑。}
\par\end{samepage}

\section{附录B｜跨册复习及长句训练}\label{ux9644ux5f55bux8de8ux518cux590dux4e60ux53caux957fux53e5ux8badux7ec3}

\begin{samepage}
\textbf{R01}　A resilient forest can recover from some disturbances, but its recovery depends on local conditions.\\
\zhline{有恢复力的森林能够从某些扰动中恢复，但恢复情况取决于当地条件。}
\par\end{samepage}

\begin{samepage}
\textbf{R02}　Subsequent reports questioned whether the initial improvement would last.\\
\zhline{后续报告质疑最初的改善能否持续。}
\par\end{samepage}

\begin{samepage}
\textbf{R03}　A government may regulate construction without having the power to prohibit every project.\\
\zhline{政府可以规范建设活动，但未必有权禁止每一个项目。}
\par\end{samepage}

\begin{samepage}
\textbf{S01}　Although the court reduced the scope of federal regulation, pollution carried by rivers could still affect communities beyond the states in which it originated.\\
\zhline{尽管法院缩小了联邦监管范围，但河流携带的污染仍可能影响污染源所在州以外的社区。}
\par\end{samepage}

\begin{samepage}
\textbf{S02}　The belief that local laws alone would be sufficient was challenged by the fact that water and the pollutants it carries do not respect administrative boundaries.\\
\zhline{认为仅凭地方法律就足够的想法，受到一个事实的挑战：水及其携带的污染物并不遵循行政边界。}
\par\end{samepage}

\begin{samepage}
\textbf{S03}　What appears to be a sensible decision for one town may have serious repercussions for people living downstream.\\
\zhline{对于一个城镇似乎合理的决定，可能给下游居民带来严重的连带影响。}
\par\end{samepage}

\begin{samepage}
\textbf{S04}　The very measures intended to protect forests from frequent fires may, over decades, increase the fuel available to a later, more destructive blaze.\\
\zhline{恰恰是那些原本旨在保护森林免遭频繁火灾的措施，经过数十年却可能增加供后续更具破坏性的大火燃烧的可燃物。}
\par\end{samepage}

\begin{samepage}
\textbf{S05}　Only after comparing modern records with evidence preserved in old trees did researchers begin to question whether fewer fires necessarily meant safer forests.\\
\zhline{只有把现代记录与保存在古树中的证据进行比较之后，研究者才开始质疑：火灾更少是否必然意味着森林更安全。}
\par\end{samepage}

\begin{samepage}
\textbf{S06}　Rather than claiming that prescribed burns can eliminate wildfire risk, the report argues that they may be one useful tool when conditions are carefully assessed.\\
\zhline{报告并未声称计划烧除可以消除野火风险，而是认为，在审慎评估条件后，计划烧除可能是一种有用的手段。}
\par\end{samepage}

\begin{samepage}
\textbf{S07}　The challenge is not simply to choose between protecting homes and restoring natural processes, but to decide which risks should be addressed immediately and which require longer-term management.\\
\zhline{难点不只是要在保护住宅与恢复自然过程之间作选择，而是要决定哪些风险需要立即处理，哪些则需要长期管理。}
\par\end{samepage}

\begin{samepage}
\textbf{S08}　That the number of recorded fires declined is evidence of successful suppression, but it does not by itself establish that the forest became safer in every respect.\\
\zhline{有记录的火灾次数下降，是积极扑火成功的证据；但仅凭这一点，还无法断定森林在各个方面都变得更加安全。}
\par\end{samepage}

\section{附录C｜开篇诊断题（含全部选项）}\label{ux9644ux5f55cux5f00ux7bc7ux8bcaux65adux9898ux542bux5168ux90e8ux9009ux9879}

\begin{samepage}
\textbf{D01}　The ruling narrowed the scope of federal oversight.\\
\zhline{这一裁决缩小了联邦监管的适用范围。}
\par\end{samepage}

\begin{samepage}
\emph{What changed?}\\
\zhline{发生了什么变化？}
\par\end{samepage}

\begin{samepage}
A. The physical size of the wetlands.\\
\zhline{湿地的实际面积。}
\par\end{samepage}

\begin{samepage}
B. The range of federal oversight.\\
\zhline{联邦监管所覆盖的范围。}
\par\end{samepage}

\begin{samepage}
C. The amount of rainfall.\\
\zhline{降雨量。}
\par\end{samepage}

\begin{samepage}
D. The number of state governments.\\
\zhline{州政府的数量。}
\par\end{samepage}

\begin{samepage}
\textbf{D02}　The proposal was at odds with the findings of earlier studies.\\
\zhline{这一提案与先前研究的发现不一致。}
\par\end{samepage}

\begin{samepage}
\emph{The relationship between the proposal and the evidence is one of \_\_\_\_\_.}\\
\zhline{提案与证据之间的关系属于哪一种？}
\par\end{samepage}

\begin{samepage}
A. agreement\\
\zhline{一致}
\par\end{samepage}

\begin{samepage}
B. uncertainty\\
\zhline{不确定}
\par\end{samepage}

\begin{samepage}
C. conflict\\
\zhline{冲突}
\par\end{samepage}

\begin{samepage}
D. repetition\\
\zhline{重复}
\par\end{samepage}

\begin{samepage}
\textbf{D03}　Measures that curb pollution do not necessarily remove every source of it.\\
\zhline{抑制污染的措施未必能消除每一种污染来源。}
\par\end{samepage}

\begin{samepage}
\emph{Which statement follows?}\\
\zhline{哪项判断符合原句？}
\par\end{samepage}

\begin{samepage}
A. The measures may reduce but not eliminate pollution.\\
\zhline{这些措施可能减少污染，却未必将其彻底消除。}
\par\end{samepage}

\begin{samepage}
B. The measures have no effect on pollution.\\
\zhline{这些措施对污染没有任何作用。}
\par\end{samepage}

\begin{samepage}
C. All pollution sources have disappeared.\\
\zhline{所有污染源均已消失。}
\par\end{samepage}

\begin{samepage}
D. The measures always increase pollution.\\
\zhline{这些措施总会增加污染。}
\par\end{samepage}

\begin{samepage}
\textbf{D04}　A policy meant to reduce short-term harm may create risks that emerge only later.\\
\zhline{一项旨在减少短期损害的政策，可能产生直到后来才显现的风险。}
\par\end{samepage}

\begin{samepage}
\emph{What does the sentence emphasize?}\\
\zhline{这句话强调什么？}
\par\end{samepage}

\begin{samepage}
A. The absence of policy benefits.\\
\zhline{政策毫无益处。}
\par\end{samepage}

\begin{samepage}
B. The certainty of policy failure.\\
\zhline{政策必然失败。}
\par\end{samepage}

\begin{samepage}
C. The difficulty of passing new laws.\\
\zhline{通过新法律的困难。}
\par\end{samepage}

\begin{samepage}
D. The possible tension between immediate and later effects.\\
\zhline{即时效果与后续影响之间可能存在张力。}
\par\end{samepage}

\section{附录D｜原创阅读A、B全文及题目}\label{ux9644ux5f55dux539fux521bux9605ux8bfbabux5168ux6587ux53caux9898ux76ee}

\subsection{Reading A｜Where Flood Protection Begins}\label{reading-awhere-flood-protection-begins-1}

\begin{samepage}
\textbf{P1}　For years, a coastal town treated heavy rain as a problem of drainage. When its streets flooded, officials widened pipes and cleared blocked channels. These measures offered visible relief, and residents could see where their taxes had gone. But as storms became more intense, the town discovered a less convenient fact: some of the water that filled its streets had arrived from developments far beyond its borders.\\
\zhline{多年来，一座沿海小镇一直把暴雨视为排水问题。街道被淹时，官员们就拓宽管道、清理堵塞的水渠。这些措施带来了看得见的缓解效果，居民也能看到税款花在了哪里。然而，随着暴雨愈发猛烈，小镇发现了一个不那么方便面对的事实：淹没街道的部分雨水来自辖区之外的开发区。}
\par\end{samepage}

\begin{samepage}
\textbf{P2}　The surrounding region had lost patches of grassland and wetland to roads, warehouses and houses. Such development did not necessarily cause flooding at any one site. Yet by reducing the ground's ability to absorb rain, it allowed more runoff to move quickly into streams that passed through the town. The effect of a planning decision was therefore not confined to the jurisdiction in which that decision had been made.\\
\zhline{周边地区的一些草地和湿地被道路、仓库与住宅取代。此类开发未必会使每一处开发地点发生洪涝，但由于降低了地面吸收雨水的能力，它使更多地表径流迅速进入流经小镇的溪流。因此，一项规划决策的影响并不局限于作出该决策的行政区域。}
\par\end{samepage}

\begin{samepage}
\textbf{P3}　Some town councillors called on the regional government to introduce rules limiting construction on remaining wetlands. Builders argued that the proposal was at odds with the demand for affordable housing. They also questioned whether the town had the authority to influence decisions made elsewhere. Both concerns deserved attention. Housing was genuinely needed, and the legal scope of local government could not simply be expanded because a problem was urgent.\\
\zhline{部分镇议员呼吁地区政府出台规则，限制在剩余湿地上进行建设。开发商认为，该提案与对可负担住房的需求相冲突。他们还质疑，小镇是否有权影响其他地区作出的决策。这两方面的担忧都值得关注。人们确实需要住房，而且不能因为问题紧迫，就随意扩大地方政府的法定权限。}
\par\end{samepage}

\begin{samepage}
\textbf{P4}　However, treating the dispute as a choice between building homes and protecting nature missed the larger issue. The costs of a project were not always borne by those who benefited from it. An upstream development could increase flood risks downstream, while the town paying for larger drains had little control over the land-use decisions that made those drains necessary. A regional review would not settle every disagreement, but it could make those connections more visible.\\
\zhline{然而，把争议看成``建住房还是保护自然''的二选一，忽略了更大的问题。项目成本并不总由项目受益者承担。上游开发可能增加下游的洪涝风险，而花钱扩建排水设施的小镇，却很难控制那些使扩建变得必要的土地利用决策。地区层面的审查无法解决所有分歧，却能够让这些联系更加清晰。}
\par\end{samepage}

\begin{samepage}
\textbf{P5}　The lesson is less about finding a perfect rule than about deciding at what level a problem should be addressed. Pipes remain useful; so do carefully planned housing projects. But where water crosses administrative boundaries, the effects of development should be assessed across those boundaries as well. Ignoring that connection may be the most shortsighted policy of all.\\
\zhline{这里的启示，与其说在于找到一条完美规则，不如说在于判断应当在哪个层面处理问题。排水管道仍然有用，经过审慎规划的住房项目也是如此。但是，既然水能跨越行政边界，开发活动的影响也应当跨越这些边界来评估。忽视这种联系，或许才是最缺乏长远眼光的政策。}
\par\end{samepage}

\subsection{Reading A｜题目与选项对照}\label{reading-aux9898ux76eeux4e0eux9009ux9879ux5bf9ux7167}

\Needspace{10\baselineskip}
\begin{samepage}
\textbf{A1.}　What did the town discover as storms grew stronger?\\
\zhline{随着暴雨增强，小镇发现了什么？}
\par\end{samepage}

\begin{samepage}
A. Most flooding resulted from faulty pipes alone.\\
\zhline{大多数洪涝仅由排水管故障导致。}
\par\end{samepage}

\begin{samepage}
B. Some floodwater originated outside the town.\\
\zhline{部分洪水来自小镇辖区之外。}
\par\end{samepage}

\begin{samepage}
C. Wider streets could prevent damage more effectively.\\
\zhline{拓宽道路能够更有效地防止损失。}
\par\end{samepage}

\begin{samepage}
D. Residents had stopped supporting drainage projects.\\
\zhline{居民已经不再支持排水工程。}
\par\end{samepage}

\Needspace{10\baselineskip}
\begin{samepage}
\textbf{A2.}　The effect of surrounding development is mainly explained by its impact on \_\_\_\_\_.\\
\zhline{周边开发的影响主要通过它对什么的作用得到解释？}
\par\end{samepage}

\begin{samepage}
A. the speed and amount of runoff\\
\zhline{径流的速度和数量}
\par\end{samepage}

\begin{samepage}
B. the level of public housing demand\\
\zhline{公共住房的需求水平}
\par\end{samepage}

\begin{samepage}
C. the quality of local tax services\\
\zhline{地方税收服务的质量}
\par\end{samepage}

\begin{samepage}
D. the number of regional officials\\
\zhline{地区官员的数量}
\par\end{samepage}

\Needspace{10\baselineskip}
\begin{samepage}
\textbf{A3.}　Why does the author acknowledge the builders' objections?\\
\zhline{作者为什么承认开发商的反对意见值得注意？}
\par\end{samepage}

\begin{samepage}
A. To argue that wetlands provide no useful services.\\
\zhline{为了论证湿地没有有益作用。}
\par\end{samepage}

\begin{samepage}
B. To prove that the town has unlimited regulatory power.\\
\zhline{为了证明小镇拥有无限监管权。}
\par\end{samepage}

\begin{samepage}
C. To recognize competing needs and legal limits.\\
\zhline{为了承认彼此竞争的需求与法律界限。}
\par\end{samepage}

\begin{samepage}
D. To recommend abandoning all regional cooperation.\\
\zhline{为了建议放弃全部地区合作。}
\par\end{samepage}

\Needspace{10\baselineskip}
\begin{samepage}
\textbf{A4.}　The problem with assessing each development separately is that \_\_\_\_\_.\\
\zhline{单独评估每个开发项目的问题在于什么？}
\par\end{samepage}

\begin{samepage}
A. the benefits of new homes cannot be measured\\
\zhline{新住房带来的收益无法衡量}
\par\end{samepage}

\begin{samepage}
B. the local government cannot collect rainfall data\\
\zhline{地方政府无法收集降雨数据}
\par\end{samepage}

\begin{samepage}
C. some projects have no direct environmental effects\\
\zhline{有些项目没有直接环境影响}
\par\end{samepage}

\begin{samepage}
D. benefits and costs may fall on different communities\\
\zhline{受益者与承担成本者可能分属不同社区}
\par\end{samepage}

\Needspace{10\baselineskip}
\begin{samepage}
\textbf{A5.}　Which of the following best expresses the author's view?\\
\zhline{哪项最能表达作者的观点？}
\par\end{samepage}

\begin{samepage}
A. Problems that cross boundaries require broader assessment.\\
\zhline{跨越边界的问题需要更大范围的评估。}
\par\end{samepage}

\begin{samepage}
B. Environmental rules should always replace housing policies.\\
\zhline{环保规则应当始终取代住房政策。}
\par\end{samepage}

\begin{samepage}
C. Local drainage improvements are now unnecessary.\\
\zhline{地方排水改善已经不再必要。}
\par\end{samepage}

\begin{samepage}
D. Any regional review will settle every disagreement.\\
\zhline{任何地区审查都能解决全部争议。}
\par\end{samepage}

\subsection{Reading B｜When Prevention Has a Cost}\label{reading-bwhen-prevention-has-a-cost-1}

\begin{samepage}
\textbf{P1}　The residents of a mountain county once judged its fire service by a simple measure: the fewer fires, the better. Crews worked quickly to extinguish any blaze near roads, houses or public campsites. For decades, the policy seemed successful. Property losses were limited, and officials could point to a steady decline in the number of recorded fires. Few people questioned a strategy that appeared to offer such clear protection.\\
\zhline{某山区县的居民过去用一项简单指标评价消防部门：火灾越少越好。队员迅速扑灭道路、住宅或公共营地附近的任何火情。几十年来，这项政策似乎十分成功。财产损失有限，官员还可以指出有记录的火灾数量持续下降。很少有人质疑一项看起来保护效果如此明确的策略。}
\par\end{samepage}

\begin{samepage}
\textbf{P2}　Then researchers began to compare recent records with evidence preserved in old trees. They found that small fires had once passed through certain parts of the forest far more often than was commonly assumed. Many of these blazes moved slowly and petered out after consuming dry leaves and low vegetation. They did not spare every tree, but they reduced the fuel available to later fires and opened space for new growth. Preventing all such fires was not necessarily the same as making the forest safer.\\
\zhline{后来，研究者开始把近期记录与保存在古树中的证据相比较。他们发现，过去某些林区发生小规模火灾的频率远高于人们通常的想象。许多火势蔓延缓慢，燃尽枯叶和低矮植被后便逐渐熄灭。它们并未让每棵树都免受伤害，却减少了可供后续火灾燃烧的燃料，还为新生植被腾出了空间。阻止所有此类火灾，并不必然等于使森林更安全。}
\par\end{samepage}

\begin{samepage}
\textbf{P3}　The finding prompted an uncomfortable question. If successful suppression allowed dead branches to accumulate, could past protection have increased some future risks? The answer was not straightforward. Warmer, drier conditions could also allow fires to spread more rapidly, and a single season of extreme weather might overwhelm even well-managed land. Nevertheless, the evidence suggested that judging policy by fire counts alone was shortsighted. Frequency, severity and potential damage were different matters.\\
\zhline{这一发现引出了一个令人不安的问题：如果成功的扑火让枯枝持续堆积，那么过去的保护是否可能增加某些未来风险？答案并不简单。更温暖、更干燥的条件也可能使火势蔓延得更快，而一个极端天气季节甚至可能使管理良好的林地也难以承受。不过，证据表明，仅凭火灾发生次数评价政策是缺乏长远考虑的。频率、严重程度与潜在损失是不同的问题。}
\par\end{samepage}

\begin{samepage}
\textbf{P4}　Officials proposed a limited programme of prescribed burns in carefully chosen locations. Small areas would be treated when wind, humidity and staffing conditions were suitable. Residents welcomed the promise of lower fuel loads but worried about smoke, safety and the possibility of a planned fire escaping control. Those concerns were not merely obstacles to public understanding: they were practical costs that the programme had to address. Officials therefore agreed to publish the conditions under which burns could proceed and to postpone operations when those conditions were not met.\\
\zhline{官员们提出在精心选择的地点开展有限范围的计划烧除。只有在风力、湿度和人员条件适宜时，才处理小片区域。居民欢迎降低可燃物负荷的前景，却担心烟雾、安全，以及计划烧除失控的可能。这些担忧不仅仅是公众理解上的障碍，而是该计划必须处理的实际成本。因此，官员同意公开可以实施烧除的条件，并在条件不符合时推迟行动。}
\par\end{samepage}

\begin{samepage}
\textbf{P5}　The county did not exchange one simple rule for another. It began to measure success by several outcomes rather than one, including the protection of homes, the health of the forest and the risks created by each management choice. That may produce less impressive numbers in the short term. But it offers a more realistic way to balance immediate safety with the ability to withstand future fires.\\
\zhline{该县并没有用另一条简单规则替换原来的规则。它开始从多个结果而非单一指标衡量成效，包括住宅保护、森林健康和各种管理选择带来的风险。这在短期内可能产生不那么好看的数字，但它为平衡眼前安全与抵御未来火灾的能力提供了更现实的办法。}
\par\end{samepage}

\subsection{Reading B｜题目与选项对照}\label{reading-bux9898ux76eeux4e0eux9009ux9879ux5bf9ux7167}

\Needspace{10\baselineskip}
\begin{samepage}
\textbf{B1.}　Why had the county's fire policy seemed successful?\\
\zhline{为什么该县的火灾政策过去似乎很成功？}
\par\end{samepage}

\begin{samepage}
A. It prevented all changes in forest conditions.\\
\zhline{它防止了森林条件发生任何变化。}
\par\end{samepage}

\begin{samepage}
B. It eliminated the need for professional fire crews.\\
\zhline{它消除了对专业消防队伍的需求。}
\par\end{samepage}

\begin{samepage}
C. Recorded fires fell and property losses were limited.\\
\zhline{记录中的火灾减少，财产损失得到控制。}
\par\end{samepage}

\begin{samepage}
D. It made prescribed burns widely accepted.\\
\zhline{它使计划烧除得到广泛接受。}
\par\end{samepage}

\Needspace{10\baselineskip}
\begin{samepage}
\textbf{B2.}　The researchers' comparison of old and recent records indicated that \_\_\_\_\_.\\
\zhline{研究者对新旧记录的比较表明什么？}
\par\end{samepage}

\begin{samepage}
A. many small fires once contributed to forest renewal\\
\zhline{过去许多小火有助于森林更新}
\par\end{samepage}

\begin{samepage}
B. early forests experienced no destructive fires\\
\zhline{早期森林从不发生破坏性火灾}
\par\end{samepage}

\begin{samepage}
C. tree-ring evidence was too unreliable to be useful\\
\zhline{树轮证据太不可靠，毫无用处}
\par\end{samepage}

\begin{samepage}
D. fire activity remained unchanged over time\\
\zhline{火灾活动始终没有变化}
\par\end{samepage}

\Needspace{10\baselineskip}
\begin{samepage}
\textbf{B3.}　What does the author suggest about the number of fires?\\
\zhline{作者对火灾次数这一指标有何看法？}
\par\end{samepage}

\begin{samepage}
A. It provides a complete measure of environmental safety.\\
\zhline{它完整地衡量了环境安全。}
\par\end{samepage}

\begin{samepage}
B. It is irrelevant to evaluating fire policy.\\
\zhline{它与火灾政策评估完全无关。}
\par\end{samepage}

\begin{samepage}
C. It should replace all measures of property damage.\\
\zhline{它应取代所有财产损失指标。}
\par\end{samepage}

\begin{samepage}
D. It cannot alone show how safe a forest has become.\\
\zhline{它不能单独说明森林已经变得多安全。}
\par\end{samepage}

\Needspace{10\baselineskip}
\begin{samepage}
\textbf{B4.}　Residents' concerns about prescribed burns are described as \_\_\_\_\_.\\
\zhline{文中如何描述居民对计划烧除的担忧？}
\par\end{samepage}

\begin{samepage}
A. proof that the entire plan should be cancelled\\
\zhline{整个计划必须取消的证明}
\par\end{samepage}

\begin{samepage}
B. practical issues to be managed in implementation\\
\zhline{实施过程中需要处理的实际问题}
\par\end{samepage}

\begin{samepage}
C. misunderstandings caused only by poor education\\
\zhline{纯粹由教育不足造成的误解}
\par\end{samepage}

\begin{samepage}
D. evidence that suppression is always preferable\\
\zhline{证明积极扑火始终更好的证据}
\par\end{samepage}

\Needspace{10\baselineskip}
\begin{samepage}
\textbf{B5.}　The main lesson of the text is that fire management should \_\_\_\_\_.\\
\zhline{本文的主要启示是火灾管理应当怎样？}
\par\end{samepage}

\begin{samepage}
A. focus solely on reducing annual fire counts\\
\zhline{只关注降低每年火灾次数}
\par\end{samepage}

\begin{samepage}
B. ignore economic costs in favour of natural processes\\
\zhline{为了自然过程而忽略经济成本}
\par\end{samepage}

\begin{samepage}
C. weigh immediate safety and long-term resilience together\\
\zhline{同时权衡眼前安全与长期恢复力}
\par\end{samepage}

\begin{samepage}
D. allow every fire to burn without restriction\\
\zhline{不加限制地允许每场火燃烧}
\par\end{samepage}

\section{附录E｜专项练习逐句对照}\label{ux9644ux5f55eux4e13ux9879ux7ec3ux4e60ux9010ux53e5ux5bf9ux7167}

\textbf{选词填空}

\begin{samepage}
\textbf{V01}　The committee will \_\_\_\_\_\_ information from several agencies before reaching a decision.\\
\zhline{委员会将在作出决定前汇集几个机构的信息。}
\par\end{samepage}

\begin{samepage}
\textbf{V02}　The river runs through three separate \_\_\_\_\_\_, making joint management difficult.\\
\zhline{这条河流经三个独立的管辖区域，使联合治理面临困难。}
\par\end{samepage}

\begin{samepage}
\textbf{V03}　Over the past year, a lack of investment has gradually \_\_\_\_\_\_ confidence in the service.\\
\zhline{过去一年里，投资不足逐渐削弱了人们对这项服务的信心。}
\par\end{samepage}

\begin{samepage}
\textbf{V04}　Heavy rain carried polluted \_\_\_\_\_\_ from farms into the lake.\\
\zhline{暴雨把来自农田的受污染径流带入湖泊。}
\par\end{samepage}

\begin{samepage}
\textbf{V05}　The study lies outside the \_\_\_\_\_\_ of the current review.\\
\zhline{这一研究超出了本次审查的范围。}
\par\end{samepage}

\begin{samepage}
\textbf{V06}　Yesterday, the editorial \_\_\_\_\_\_ officials to consider the interests of nearby residents.\\
\zhline{昨天，这篇社论呼吁官员考虑附近居民的利益。}
\par\end{samepage}

\begin{samepage}
\textbf{V07}　The budget \_\_\_\_\_\_ led to a delay in repairing damaged bridges.\\
\zhline{预算缺口导致受损桥梁的维修延迟。}
\par\end{samepage}

\begin{samepage}
\textbf{V08}　The new law gave the agency the \_\_\_\_\_\_ to inspect unsafe buildings.\\
\zhline{新法律赋予该机构检查不安全建筑的权限。}
\par\end{samepage}

\begin{samepage}
\textbf{V09}　The small fire \_\_\_\_\_\_ when it reached a patch of wet ground.\\
\zhline{小火蔓延到一片潮湿的地面后逐渐熄灭。}
\par\end{samepage}

\begin{samepage}
\textbf{V10}　The plan was criticized as \_\_\_\_\_\_ because it ignored maintenance costs.\\
\zhline{该计划因忽视维护成本而被批评为缺乏长远眼光。}
\par\end{samepage}

\begin{samepage}
\textbf{V11}　The fire service carried out a \_\_\_\_\_\_ burn under suitable weather conditions.\\
\zhline{消防机构在适宜天气条件下开展了一次计划烧除。}
\par\end{samepage}

\begin{samepage}
\textbf{V12}　The decision could have unexpected \_\_\_\_\_\_ for towns downstream.\\
\zhline{这一决定可能给下游城镇带来意料之外的连带影响。}
\par\end{samepage}

\begin{samepage}
\textbf{V13}　New measures were introduced to \_\_\_\_\_\_ the spread of the disease.\\
\zhline{新措施被用来抑制这种疾病的传播。}
\par\end{samepage}

\begin{samepage}
\textbf{V14}　The two departments appeared to be \_\_\_\_\_\_ one another.\\
\zhline{两个部门似乎存在分歧。}
\par\end{samepage}

\begin{samepage}
\textbf{V15}　Several years of fire \_\_\_\_\_\_ increased the amount of dry wood on the forest floor.\\
\zhline{多年的积极扑火增加了林地地表的枯木数量。}
\par\end{samepage}

\begin{samepage}
\textbf{V16}　An agreement may help \_\_\_\_\_\_ the political divisions that hinder cooperation.\\
\zhline{协议可能有助于跨越阻碍合作的政治分歧。}
\par\end{samepage}

\textbf{词义判断}

\begin{samepage}
\textbf{M01}　The law limits the agency's authority, not the river's flow.\\
\zhline{法律限制的是机构的权限，而不是河水的流量。}
\par\end{samepage}

\begin{samepage}
A. legal power or permission\\
\zhline{法律赋予的权力或权限}
\par\end{samepage}

\begin{samepage}
B. an expert respected for knowledge\\
\zhline{因知识受到尊重的专家}
\par\end{samepage}

\begin{samepage}
\textbf{M02}　The court changed the scope of regulation, but the ecological consequences may unfold later.\\
\zhline{法院改变了监管范围，但生态后果可能在以后才逐渐显现。}
\par\end{samepage}

\begin{samepage}
A. all environmental rules have been removed\\
\zhline{全部环境规则被取消}
\par\end{samepage}

\begin{samepage}
B. the range of regulation has changed\\
\zhline{监管的适用范围发生变化}
\par\end{samepage}

\begin{samepage}
\textbf{M03}　A prescribed burn is not the same as allowing any wildfire to spread unchecked.\\
\zhline{计划烧除并不等于任由任何野火不受控制地蔓延。}
\par\end{samepage}

\begin{samepage}
A. a carefully planned fire\\
\zhline{经过审慎计划的烧除}
\par\end{samepage}

\begin{samepage}
B. any fire allowed to spread freely\\
\zhline{任由任何火自由蔓延}
\par\end{samepage}

\begin{samepage}
\textbf{M04}　A decline in the number of fires does not establish that future losses will decline.\\
\zhline{火灾次数减少，并不能证明未来损失也会减少。}
\par\end{samepage}

\begin{samepage}
A. fewer fires prove all risks are lower\\
\zhline{火灾较少就证明所有风险都降低}
\par\end{samepage}

\begin{samepage}
B. fewer fires do not prove every risk is lower\\
\zhline{火灾较少不能证明所有风险都降低}
\par\end{samepage}

\begin{samepage}
\textbf{M05}　The findings were a testament to the benefits of careful monitoring.\\
\zhline{这些发现证明了审慎监测的益处。}
\par\end{samepage}

\begin{samepage}
A. evidence or a clear sign of something\\
\zhline{某事物的证据或明显体现}
\par\end{samepage}

\begin{samepage}
B. a legal document about inheritance\\
\zhline{关于继承的法律文件}
\par\end{samepage}

\begin{samepage}
\textbf{M06}　The policy is a double-edged sword: it reduces one risk while creating another.\\
\zhline{这项政策是一把双刃剑：它降低一种风险，同时又产生另一种风险。}
\par\end{samepage}

\begin{samepage}
A. a policy free of trade-offs\\
\zhline{毫无取舍代价的政策}
\par\end{samepage}

\begin{samepage}
B. a policy with both benefits and disadvantages\\
\zhline{同时有益处和不利之处的政策}
\par\end{samepage}

\textbf{限定与推断}

\begin{samepage}
\textbf{I01}　The wetland review may affect several counties, although only one county requested it.\\
\zhline{虽然只有一个县提出开展湿地审查，但审查可能影响多个县。}
\par\end{samepage}

\begin{samepage}
\emph{The sentence allows which inference?}\\
\zhline{该句允许哪项推断？}
\par\end{samepage}

\begin{samepage}
A. The review must be limited to the requesting county.\\
\zhline{审查必须只限于提出请求的县。}
\par\end{samepage}

\begin{samepage}
B. The review has already changed every local law.\\
\zhline{审查已经改变了所有地方法律。}
\par\end{samepage}

\begin{samepage}
C. The possible effects extend beyond the requesting county.\\
\zhline{可能的影响超出了提出请求的县。}
\par\end{samepage}

\begin{samepage}
D. No county has a legal interest in the review.\\
\zhline{没有任何县在审查中具有法律利益。}
\par\end{samepage}

\begin{samepage}
\textbf{I02}　The measure reduced flooding in the town centre but increased runoff into a nearby village.\\
\zhline{该措施减少了镇中心的洪涝，却增加了流入附近村庄的径流。}
\par\end{samepage}

\begin{samepage}
\emph{What is the best conclusion?}\\
\zhline{最恰当的结论是什么？}
\par\end{samepage}

\begin{samepage}
A. The project had different effects in different places.\\
\zhline{该项目在不同地点产生了不同影响。}
\par\end{samepage}

\begin{samepage}
B. The project eliminated flooding across the region.\\
\zhline{该项目消除了整个地区的洪涝。}
\par\end{samepage}

\begin{samepage}
C. The village received no additional runoff.\\
\zhline{村庄没有接收到更多径流。}
\par\end{samepage}

\begin{samepage}
D. The project was necessarily unlawful.\\
\zhline{该项目必然违法。}
\par\end{samepage}

\begin{samepage}
\textbf{I03}　The study suggests that a policy can succeed by one measure while failing by another.\\
\zhline{研究表明，一项政策可能按某个指标算成功，按另一指标却算失败。}
\par\end{samepage}

\begin{samepage}
\emph{The main idea is that \_\_\_\_\_.}\\
\zhline{这句话的主旨是什么？}
\par\end{samepage}

\begin{samepage}
A. one indicator always determines success\\
\zhline{单一指标总能决定成败}
\par\end{samepage}

\begin{samepage}
B. policies should never be evaluated\\
\zhline{永远不应评估政策}
\par\end{samepage}

\begin{samepage}
C. all results are equally important\\
\zhline{所有结果都同样重要}
\par\end{samepage}

\begin{samepage}
D. policy outcomes may depend on how success is measured\\
\zhline{政策结果的评价可能取决于如何衡量成功}
\par\end{samepage}

\begin{samepage}
\textbf{I04}　Prescribed burns may help reduce fuel, but they must be postponed when weather conditions are unsuitable.\\
\zhline{计划烧除可能减少可燃物，但当天气条件不适合时必须推迟实施。}
\par\end{samepage}

\begin{samepage}
\emph{Which statement is consistent with this?}\\
\zhline{哪项与原句一致？}
\par\end{samepage}

\begin{samepage}
A. Fire managers can burn whenever they choose.\\
\zhline{管理者可以随时烧除。}
\par\end{samepage}

\begin{samepage}
B. The method has potential benefits but requires limits.\\
\zhline{这种方法有潜在益处，但需要遵守限制条件。}
\par\end{samepage}

\begin{samepage}
C. Prescribed burns never produce smoke.\\
\zhline{计划烧除从不产生烟雾。}
\par\end{samepage}

\begin{samepage}
D. Poor weather makes all forest management useless.\\
\zhline{恶劣天气使所有森林管理都失去意义。}
\par\end{samepage}

\textbf{英译汉}

\begin{samepage}
\textbf{T01}　A reduction in federal authority does not automatically remove every form of protection available under state law.\\
\zhline{联邦权限的缩减，并不自动取消州法下所有可用的保护措施。}
\par\end{samepage}

\begin{samepage}
\textbf{T02}　The costs of pollution may be felt far from the place where decisions about land use were originally made.\\
\zhline{污染的代价可能在距离土地利用决策最初作出地点很远的地方显现。}
\par\end{samepage}

\begin{samepage}
\textbf{T03}　It would be shortsighted to measure a policy only by its immediate benefits when the effects of that policy may last for decades.\\
\zhline{如果一项政策的影响可能持续数十年，却只用眼前收益评价它，这将是目光短浅的。}
\par\end{samepage}

\begin{samepage}
\textbf{T04}　A small fire that clears dry vegetation may reduce the intensity of later fires without preventing them entirely.\\
\zhline{一场清除干燥植被的小火，可能降低后续火灾的强度，却未必能彻底阻止它们发生。}
\par\end{samepage}

\begin{samepage}
\textbf{T05}　The fact that a programme has lowered the number of reported fires is not sufficient evidence that it has improved safety in every respect.\\
\zhline{一项计划减少了报告中的火灾次数，并不足以证明它已经在所有方面提高了安全程度。}
\par\end{samepage}

\begin{samepage}
\textbf{T06}　Even where a planned burn is considered beneficial, the timing of the operation must depend on local weather and available personnel.\\
\zhline{即使计划烧除被认为有益，实施时间仍须取决于当地天气和可用人员。}
\par\end{samepage}

\textbf{汉译英}

\begin{samepage}
\textbf{P01}　Not every environmental issue falls within the scope of the report.\\
\zhline{用 scope of 表达：报告没有讨论所有环境问题。}
\par\end{samepage}

\begin{samepage}
\textbf{P02}　The researchers called on local authorities to share water-quality data.\\
\zhline{用 call on sb. to do 表达：研究者呼吁地方政府分享水质数据。}
\par\end{samepage}

\begin{samepage}
\textbf{P03}　Persistent underfunding may gradually erode public confidence.\\
\zhline{用 erode 与 may 表达：长期资金不足可能逐渐削弱公众信心。}
\par\end{samepage}

\begin{samepage}
\textbf{P04}　The decision appears to be at odds with the available evidence.\\
\zhline{用 at odds with 表达：该决定似乎与现有证据不一致。}
\par\end{samepage}

\begin{samepage}
\textbf{P05}　We should compare several outcomes rather than count fires alone.\\
\zhline{用 rather than 表达：应比较多项结果，而不是只计算火灾次数。}
\par\end{samepage}

\section{附录F｜跨情境阅读C全文、题干与选项}\label{ux9644ux5f55fux8de8ux60c5ux5883ux9605ux8bfbcux5168ux6587ux9898ux5e72ux4e0eux9009ux9879}

\begin{samepage}
\textbf{P1}　A university library introduced automatic renewal to reduce the number of overdue books. Unless another reader requested an item, its loan period would be extended without any action by the borrower. After one term, overdue notices had fallen sharply. Administrators described the decline as evidence that students were managing their borrowing more responsibly. The new system also saved staff time, allowing them to answer more research enquiries. On these measures, the change looked like an obvious success.\\
\zhline{某大学图书馆为减少逾期图书引入了自动续借。除非有其他读者预约某本书，否则借期会自动延长，借阅者无须操作。一个学期后，逾期通知大幅减少。管理人员将此视为学生更负责任地管理借阅的证据。新系统还节省了工作人员的时间，让他们能回答更多研究咨询。按这些指标看，这项改变显然取得了成功。}
\par\end{samepage}

\begin{samepage}
\textbf{P2}　Yet a student committee questioned what the figures actually showed. Under the old rules, a book became overdue after a fixed period. Under the new rules, the same book could remain on loan for much longer without generating a notice. The committee did not deny that fewer warnings were welcome. It argued that a change in the definition of a recorded problem should not be confused with a change in the behaviour that had previously produced it.\\
\zhline{然而，学生委员会质疑这些数字究竟说明了什么。按照旧规则，借阅超过固定期限便算逾期；按照新规则，同一本书可以被借走更久而不触发通知。委员会并不否认警告减少是好事。它认为，记录某个问题的定义发生变化，不应与过去产生这一问题的行为发生变化混为一谈。}
\par\end{samepage}

\begin{samepage}
\textbf{P3}　To examine access more directly, the library compiled data on requests, waiting times and unsuccessful searches. Some readers received popular books sooner because the request system brought those books back. Others waited longer for less popular titles that remained with borrowers who no longer needed them. These mixed findings did not establish that automatic renewal had failed. They suggested that the benefits depended partly on how easily readers could discover and request material that was not on the shelves.\\
\zhline{为更直接地考察获取图书的情况，图书馆汇集了预约、等待时间和未找到所需资料的检索记录。一些读者更快拿到了热门书，因为预约系统促使这些书被归还。另一些读者却要更久地等待冷门书，这些书还留在已不再需要它们的借阅者手中。这些有好有坏的发现并不能证明自动续借失败了。它们提示，收益部分取决于读者能否方便地发现并预约不在书架上的资料。}
\par\end{samepage}

\begin{samepage}
\textbf{P4}　The committee called on the library to retain automatic renewal while making requests more visible in its catalogue. It also proposed a reminder asking borrowers whether they still needed a book. Some staff worried that extra messages would erode the simplicity students valued. The library therefore tested the reminder with a small group before deciding whether to extend it. A useful tactic in one setting, it reasoned, might create an unnecessary burden in another.\\
\zhline{委员会呼吁图书馆保留自动续借，同时在目录中更醒目地显示预约功能。它还提议发送提醒，询问借阅者是否仍需要某本书。一些工作人员担心额外信息会削弱学生所看重的简便性。因此，图书馆先在一个小群体中试行提醒，再决定是否推广。它认为，在一种情境下有用的做法，在另一种情境下可能造成不必要的负担。}
\par\end{samepage}

\begin{samepage}
\textbf{P5}　The lesson was not that numerical targets should be abandoned. Without records, unequal access might have remained invisible. Rather, institutions needed to ask whether a reported improvement reflected better service, altered recording rules, or both. A policy could curb one inconvenience while creating another; evaluating it required following the experience of users, not simply celebrating the disappearance of an administrative warning.\\
\zhline{这里的启示并非应当放弃数字目标。没有记录，获取机会的不均等可能一直不被察觉。机构需要追问的是，报告中的改善究竟反映服务变好了、记录规则变了，还是两者都有。一项政策可能减少一种不便，同时产生另一种不便；评价政策需要跟踪使用者的体验，而不能只庆祝管理警告的消失。}
\par\end{samepage}

\Needspace{10\baselineskip}
\begin{samepage}
\textbf{C1.} What does the committee question about the decline in overdue notices?\\
\zhline{委员会对逾期通知减少提出了什么质疑？}
\par\end{samepage}

\begin{samepage}
A. It necessarily indicates more responsible borrowing.\\
\zhline{它必然表明借阅行为更加负责任。}
\par\end{samepage}

\begin{samepage}
B. It reduced the time available for research enquiries.\\
\zhline{它减少了回答研究咨询的时间。}
\par\end{samepage}

\begin{samepage}
C. It followed the introduction of automatic renewal.\\
\zhline{它发生在引入自动续借之后。}
\par\end{samepage}

\begin{samepage}
D. It was welcomed by some library users.\\
\zhline{它受到部分图书馆使用者的欢迎。}
\par\end{samepage}

\Needspace{10\baselineskip}
\begin{samepage}
\textbf{C2.} Why are less popular books mentioned in Paragraph 3?\\
\zhline{第3段为什么提到冷门书？}
\par\end{samepage}

\begin{samepage}
A. To show that all borrowers prefer popular titles.\\
\zhline{说明所有借阅者都更喜欢热门书。}
\par\end{samepage}

\begin{samepage}
B. To establish that loan periods should never be extended.\\
\zhline{证明借期永远不应延长。}
\par\end{samepage}

\begin{samepage}
C. To illustrate differences in the effects on access.\\
\zhline{说明政策对获取图书的影响存在差异。}
\par\end{samepage}

\begin{samepage}
D. To suggest that search records should be discarded.\\
\zhline{建议丢弃检索记录。}
\par\end{samepage}

\Needspace{10\baselineskip}
\begin{samepage}
\textbf{C3.} The proposed reminder was tested with a small group because \_\_\_\_\_.\\
\zhline{所提议的提醒先在小群体中试行，是因为什么？}
\par\end{samepage}

\begin{samepage}
A. staff had proved that students disliked every message\\
\zhline{工作人员已证明学生讨厌所有消息}
\par\end{samepage}

\begin{samepage}
B. the library wanted to examine a possible trade-off\\
\zhline{图书馆想考察一种可能的利弊权衡}
\par\end{samepage}

\begin{samepage}
C. the committee wanted automatic renewal cancelled\\
\zhline{委员会希望取消自动续借}
\par\end{samepage}

\begin{samepage}
D. all books had already returned to the shelves\\
\zhline{所有书都已经归架}
\par\end{samepage}

\Needspace{10\baselineskip}
\begin{samepage}
\textbf{C4.} Which claim is supported by the text?\\
\zhline{哪一项说法得到文章支持？}
\par\end{samepage}

\begin{samepage}
A. A fall in warnings proves that access has improved for everyone.\\
\zhline{警告减少证明每个人获取图书都更方便了。}
\par\end{samepage}

\begin{samepage}
B. Changing recording rules cannot affect reported performance.\\
\zhline{改变记录规则不会影响报告中的表现。}
\par\end{samepage}

\begin{samepage}
C. Mixed results make further evaluation unnecessary.\\
\zhline{结果有好有坏便无须继续评价。}
\par\end{samepage}

\begin{samepage}
D. Service improvements and changes in recording may occur together.\\
\zhline{服务改善与记录方式改变可能同时发生。}
\par\end{samepage}

\Needspace{10\baselineskip}
\begin{samepage}
\textbf{C5.} Which title best captures the author's argument?\\
\zhline{哪个标题最能概括作者的论点？}
\par\end{samepage}

\begin{samepage}
A. The Case for Ending All Library Targets\\
\zhline{取消图书馆所有目标的理由}
\par\end{samepage}

\begin{samepage}
B. Why Borrowers Should Stop Requesting Books\\
\zhline{借阅者为何应停止预约}
\par\end{samepage}

\begin{samepage}
C. When Fewer Warnings Need a Closer Look\\
\zhline{警告减少时仍须仔细审视}
\par\end{samepage}

\begin{samepage}
D. Automatic Renewal as a Complete Failure\\
\zhline{自动续借的彻底失败}
\par\end{samepage}

\section{附录G｜参考写法补充与阅读标题}\label{ux9644ux5f55gux53c2ux8003ux5199ux6cd5ux8865ux5145ux4e0eux9605ux8bfbux6807ux9898}

\begin{samepage}
\textbf{P01补充写法}　The report does not address every environmental issue; some fall outside its scope.\\
\zhline{这份报告没有涉及每一种环境问题；有些问题超出了其讨论范围。}
\par\end{samepage}

\begin{samepage}
\textbf{Reading A}　Where Flood Protection Begins\\
\zhline{防洪应从何处着手}
\par\end{samepage}

\begin{samepage}
\textbf{Reading B}　When Prevention Has a Cost\\
\zhline{当预防也有代价}
\par\end{samepage}

\begin{samepage}
\textbf{Reading C}　When Fewer Warnings Need a Closer Look\\
\zhline{警告减少时仍须仔细审视}
\par\end{samepage}

\end{document}
